% Môn: Vật lý
% Lớp: 11
% Chương: Chương II. Sóng
% Bài: L11C2 Bài 9. Sóng ngang, sóng dọc, sự truyền năng lượng của sóng cơ
% Số câu: 170
% App đọc file này trực tiếp — sửa rồi Commit trên GitHub, bấm Đồng bộ GitHub .tex

% L11C2 Bài 9. Sóng ngang, sóng dọc, sự truyền năng lượng của sóng cơ

% ===== Câu 1 =====
% ID: L11C2B9-01-DS
% Mức: TH
\dangbt{Bài toán thực tế về âm, độ cao, độ to và môi trường truyền âm}
\begin{ex}
% ID: L11C2B9-01-DS
% Mức: TH
Trong chuyên đề “Bài toán thực tế về âm, độ cao, độ to và môi trường truyền âm”, Một loa phát âm $500\,\text{Hz}$ trong không khí.
\choiceTF
{\True Tần số âm bằng tần số dao động của nguồn.}
{\True Âm trong không khí là sóng dọc.}
{\True Tăng biên độ nguồn thường làm âm to hơn.}
{Tăng biên độ làm âm cao hơn.}
\loigiai{
\begin{itemchoice}
\item (Đúng) Tần số âm bằng tần số dao động của nguồn. (Vì): Tần số âm bằng tần số dao động của nguồn. b) Đ: Âm trong không khí là sóng dọc. c) Đ: Tăng biên độ nguồn thường làm âm to hơn. d) S: Tăng biên độ làm âm cao hơn
\item (Đúng) Âm trong không khí là sóng dọc.
\item (Đúng) Tăng biên độ nguồn thường làm âm to hơn.
\item (Sai) Tăng biên độ làm âm cao hơn.
\end{itemchoice}
}
\end{ex}

% ===== Câu 2 =====
% ID: L11C2B9-01-TL
% Mức: TH
\begin{ex}
% ID: L11C2B9-01-TL
% Mức: TH
Trong chuyên đề “Bài toán thực tế về âm, độ cao, độ to và môi trường truyền âm”, Phân biệt độ cao và độ to của âm bằng hai đặc trưng vật lí tương ứng.
\choiceTF

\loigiai{
Độ cao phụ thuộc chủ yếu vào tần số; độ to liên hệ với biên độ/cường độ âm tại tai.
}
\end{ex}

% ===== Câu 3 =====
% ID: L11C2B9-01-TLN
% Mức: TH
\begin{ex}
% ID: L11C2B9-01-TLN
% Mức: TH
Trong chuyên đề “Bài toán thực tế về âm, độ cao, độ to và môi trường truyền âm”, Âm có tốc độ $340\,\text{m}$ /s và tần số $680\,\text{Hz}$. Bước sóng tính theo m bằng bao nhiêu?
\shortans{0.5}
\loigiai{
Xác định đúng dữ kiện và đơn vị, áp dụng hệ thức của dạng bài rồi tính được kết quả 0.5.
}
\end{ex}

% ===== Câu 4 =====
% ID: L11C2B9-01-TN
% Mức: NB
\begin{ex}
% ID: L11C2B9-01-TN
% Mức: NB
Trong chuyên đề “Bài toán thực tế về âm, độ cao, độ to và môi trường truyền âm”, Độ cao của âm chủ yếu phụ thuộc vào
\choice
{biên độ}
{tốc độ gió}
{khoảng cách tới nguồn}
{\True tần số}
\loigiai{
Đáp án D. tần số. Đối chiếu định nghĩa, hệ thức hoặc dữ kiện của bài toán cho kết luận trên.
}
\end{ex}

% ===== Câu 5 =====
% ID: L11C2B9-02-DS
% Mức: TH
\begin{ex}
% ID: L11C2B9-02-DS
% Mức: TH
Trong chuyên đề “Bài toán thực tế về âm, độ cao, độ to và môi trường truyền âm”, Xét hạ âm, âm nghe được và siêu âm.
\choiceTF
{\True Hạ âm có tần số dưới ngưỡng nghe thấp.}
{\True Siêu âm có tần số lớn hơn 20 kHz theo quy ước thường dùng.}
{\True Tai người bình thường chỉ nghe được một khoảng tần số hữu hạn.}
{Mọi động vật có cùng giới hạn nghe như người.}
\loigiai{
\begin{itemchoice}
\item (Đúng) Hạ âm có tần số dưới ngưỡng nghe thấp. (Vì): Hạ âm có tần số dưới ngưỡng nghe thấp. b) Đ: Siêu âm có tần số lớn hơn 20 kHz theo quy ước thường dùng. c) Đ: Tai người bình thường chỉ nghe được một khoảng tần số hữu hạn. d) S: Mọi động vật có cùng giới hạn nghe như người
\item (Đúng) Siêu âm có tần số lớn hơn 20 kHz theo quy ước thường dùng.
\item (Đúng) Tai người bình thường chỉ nghe được một khoảng tần số hữu hạn.
\item (Sai) Mọi động vật có cùng giới hạn nghe như người.
\end{itemchoice}
}
\end{ex}

% ===== Câu 6 =====
% ID: L11C2B9-02-TL
% Mức: TH
\begin{ex}
% ID: L11C2B9-02-TL
% Mức: TH
Trong chuyên đề “Bài toán thực tế về âm, độ cao, độ to và môi trường truyền âm”, Giải thích vì sao cùng một nốt nhạc truyền từ không khí vào nước vẫn có tần số không đổi nhưng bước sóng thay đổi.
\choiceTF

\loigiai{
Tần số do nguồn quyết định và liên tục qua mặt phân cách; tốc độ phụ thuộc môi trường nên theo $\lambda=v/f$, bước sóng thay đổi.
}
\end{ex}

% ===== Câu 7 =====
% ID: L11C2B9-02-TLN
% Mức: TH
\begin{ex}
% ID: L11C2B9-02-TLN
% Mức: TH
Trong chuyên đề “Bài toán thực tế về âm, độ cao, độ to và môi trường truyền âm”, Một âm có chu kì $0,002\,\text{s}$. Tần số tính theo Hz bằng bao nhiêu?
\shortans{500}
\loigiai{
Xác định đúng dữ kiện và đơn vị, áp dụng hệ thức của dạng bài rồi tính được kết quả 500.
}
\end{ex}

% ===== Câu 8 =====
% ID: L11C2B9-02-TN
% Mức: NB
\begin{ex}
% ID: L11C2B9-02-TN
% Mức: NB
Trong chuyên đề “Bài toán thực tế về âm, độ cao, độ to và môi trường truyền âm”, Độ to của âm liên hệ trực tiếp nhất với
\choice
{\True biên độ dao động âm}
{chu kì âm}
{bước sóng duy nhất}
{pha ban đầu}
\loigiai{
Đáp án A. biên độ dao động âm. Đối chiếu định nghĩa, hệ thức hoặc dữ kiện của bài toán cho kết luận trên.
}
\end{ex}

% ===== Câu 9 =====
% ID: L11C2B9-03-DS
% Mức: TH
\begin{ex}
% ID: L11C2B9-03-DS
% Mức: TH
Trong chuyên đề “Bài toán thực tế về âm, độ cao, độ to và môi trường truyền âm”, Âm truyền từ không khí vào nước.
\choiceTF
{\True Tần số do nguồn quyết định nên không đổi.}
{\True Tốc độ truyền có thể thay đổi.}
{\True Bước sóng thay đổi khi tốc độ thay đổi.}
{Âm truyền được trong chân không giữa hai môi trường.}
\loigiai{
\begin{itemchoice}
\item (Đúng) Tần số do nguồn quyết định nên không đổi. (Vì): Tần số do nguồn quyết định nên không đổi. b) Đ: Tốc độ truyền có thể thay đổi. c) Đ: Bước sóng thay đổi khi tốc độ thay đổi. d) S: Âm truyền được trong chân không giữa hai môi trường
\item (Đúng) Tốc độ truyền có thể thay đổi.
\item (Đúng) Bước sóng thay đổi khi tốc độ thay đổi.
\item (Sai) Âm truyền được trong chân không giữa hai môi trường.
\end{itemchoice}
}
\end{ex}

% ===== Câu 10 =====
% ID: L11C2B9-03-TL
% Mức: VD
\begin{ex}
% ID: L11C2B9-03-TL
% Mức: VD
Trong chuyên đề “Bài toán thực tế về âm, độ cao, độ to và môi trường truyền âm”, Một âm $850\,\text{Hz}$ truyền trong không khí với tốc độ $340\,\text{m}$ /s. Tính bước sóng và cho biết âm thuộc vùng nghe được hay không.
\choiceTF

\loigiai{
$\lambda=340/850=0,4$ m. Tần số $850\,\text{Hz}$ thuộc vùng âm nghe được của người.
}
\end{ex}

% ===== Câu 11 =====
% ID: L11C2B9-03-TLN
% Mức: TH
\begin{ex}
% ID: L11C2B9-03-TLN
% Mức: TH
Trong chuyên đề “Bài toán thực tế về âm, độ cao, độ to và môi trường truyền âm”, Siêu âm tần số 40 kHz truyền với tốc độ $320\,\text{m}$ /s. Bước sóng tính theo mm bằng bao nhiêu?
\shortans{8}
\loigiai{
Xác định đúng dữ kiện và đơn vị, áp dụng hệ thức của dạng bài rồi tính được kết quả 8.
}
\end{ex}

% ===== Câu 12 =====
% ID: L11C2B9-03-TN
% Mức: TH
\begin{ex}
% ID: L11C2B9-03-TN
% Mức: TH
Trong chuyên đề “Bài toán thực tế về âm, độ cao, độ to và môi trường truyền âm”, Âm có tần số 25 kHz được gọi là
\choice
{sóng vô tuyến}
{\True siêu âm}
{hạ âm}
{âm nghe được}
\loigiai{
Đáp án B. siêu âm. Đối chiếu định nghĩa, hệ thức hoặc dữ kiện của bài toán cho kết luận trên.
}
\end{ex}

% ===== Câu 13 =====
% ID: L11C2B9-04-DS
% Mức: VD
\begin{ex}
% ID: L11C2B9-04-DS
% Mức: VD
Trong chuyên đề “Bài toán thực tế về âm, độ cao, độ to và môi trường truyền âm”, Xét các nhận định tổng hợp thuộc dạng “Bài toán thực tế về âm, độ cao, độ to và môi trường truyền âm”.
\choiceTF
{\True tần số.}
{\True biên độ dao động âm.}
{\True siêu âm.}
{biên độ.}
\loigiai{
\begin{itemchoice}
\item (Đúng) tần số. (Vì): tần số. b) Đ: biên độ dao động âm. c) Đ: siêu âm. d) S: biên độ
\item (Đúng) biên độ dao động âm.
\item (Đúng) siêu âm.
\item (Sai) biên độ.
\end{itemchoice}
}
\end{ex}

% ===== Câu 14 =====
% ID: L11C2B9-04-TL
% Mức: VD
\begin{ex}
% ID: L11C2B9-04-TL
% Mức: VD
Trong chuyên đề “Bài toán thực tế về âm, độ cao, độ to và môi trường truyền âm”, Phân tích các bước lập luận và chỉ ra điều kiện áp dụng trong bài toán: Phân biệt độ cao và độ to của âm bằng hai đặc trưng vật lí tương ứng.
\choiceTF

\loigiai{
Độ cao phụ thuộc chủ yếu vào tần số; độ to liên hệ với biên độ/cường độ âm tại tai. Cần nêu rõ điều kiện áp dụng, đổi đơn vị thống nhất và kiểm tra ý nghĩa vật lí của kết quả.
}
\end{ex}

% ===== Câu 15 =====
% ID: L11C2B9-04-TLN
% Mức: VD
\begin{ex}
% ID: L11C2B9-04-TLN
% Mức: VD
Trong chuyên đề “Bài toán thực tế về âm, độ cao, độ to và môi trường truyền âm”, Từ kết quả của tình huống sau, nếu đại lượng cần tìm được nhân với hệ số 2 thì giá trị mới bằng bao nhiêu? Âm có tốc độ $340\,\text{m}$ /s và tần số $680\,\text{Hz}$. Bước sóng tính theo m bằng bao nhiêu?
\shortans{1}
\loigiai{
Xác định đúng dữ kiện và đơn vị, áp dụng hệ thức của dạng bài rồi tính được kết quả 1.
}
\end{ex}

% ===== Câu 16 =====
% ID: L11C2B9-04-TN
% Mức: TH
\begin{ex}
% ID: L11C2B9-04-TN
% Mức: TH
Trong chuyên đề “Bài toán thực tế về âm, độ cao, độ to và môi trường truyền âm”, Âm có tốc độ $340\,\text{m}$ /s và tần số $680\,\text{Hz}$. Bước sóng tính theo m bằng bao nhiêu? Chọn giá trị đúng.
\choice
{0.25}
{01/05/2026}
{\True 0.5}
{1}
\loigiai{
Đáp án C. 0.5. Đối chiếu định nghĩa, hệ thức hoặc dữ kiện của bài toán cho kết luận trên.
}
\end{ex}

% ===== Câu 17 =====
% ID: L11C2B9-05-DS
% Mức: VD
\begin{ex}
% ID: L11C2B9-05-DS
% Mức: VD
Trong chuyên đề “Bài toán thực tế về âm, độ cao, độ to và môi trường truyền âm”, Đánh giá cách xử lí các dữ kiện định lượng của dạng “Bài toán thực tế về âm, độ cao, độ to và môi trường truyền âm”.
\choiceTF
{\True Kết quả của bài toán thứ nhất là 0.5.}
{\True Kết quả của bài toán thứ hai là 500.}
{\True Kết quả của bài toán thứ ba là 8.}
{Có thể bỏ qua hoàn toàn đơn vị và điều kiện vật lí khi tính toán.}
\loigiai{
\begin{itemchoice}
\item (Đúng) Kết quả của bài toán thứ nhất là 0.5. (Vì): Kết quả của bài toán thứ nhất là 0.5. b) Đ: Kết quả của bài toán thứ hai là 500. c) Đ: Kết quả của bài toán thứ ba là 8. d) S: Có thể bỏ qua hoàn toàn đơn vị và điều kiện vật lí khi tính toán
\item (Đúng) Kết quả của bài toán thứ hai là 500.
\item (Đúng) Kết quả của bài toán thứ ba là 8.
\item (Sai) Có thể bỏ qua hoàn toàn đơn vị và điều kiện vật lí khi tính toán.
\end{itemchoice}
}
\end{ex}

% ===== Câu 18 =====
% ID: L11C2B9-05-TL
% Mức: VD
\begin{ex}
% ID: L11C2B9-05-TL
% Mức: VD
Trong chuyên đề “Bài toán thực tế về âm, độ cao, độ to và môi trường truyền âm”, Một học sinh giải bài sau nhưng bỏ qua điều kiện vật lí. Hãy sửa cách làm và trình bày lời giải đúng: Giải thích vì sao cùng một nốt nhạc truyền từ không khí vào nước vẫn có tần số không đổi nhưng bước sóng thay đổi.
\choiceTF

\loigiai{
Cách giải đúng: Tần số do nguồn quyết định và liên tục qua mặt phân cách; tốc độ phụ thuộc môi trường nên theo $\lambda=v/f$, bước sóng thay đổi. Sai lầm cần tránh là dùng hệ thức khi chưa xác định điều kiện biên, môi trường hoặc đơn vị.
}
\end{ex}

% ===== Câu 19 =====
% ID: L11C2B9-05-TLN
% Mức: VD
\begin{ex}
% ID: L11C2B9-05-TLN
% Mức: VD
Trong chuyên đề “Bài toán thực tế về âm, độ cao, độ to và môi trường truyền âm”, Từ kết quả của tình huống sau, nếu đại lượng cần tìm được nhân với hệ số 0.5 thì giá trị mới bằng bao nhiêu? Một âm có chu kì $0,002\,\text{s}$. Tần số tính theo Hz bằng bao nhiêu?
\shortans{250}
\loigiai{
Xác định đúng dữ kiện và đơn vị, áp dụng hệ thức của dạng bài rồi tính được kết quả 250.
}
\end{ex}

% ===== Câu 20 =====
% ID: L11C2B9-05-TN
% Mức: TH
\begin{ex}
% ID: L11C2B9-05-TN
% Mức: TH
Trong chuyên đề “Bài toán thực tế về âm, độ cao, độ to và môi trường truyền âm”, Một âm có chu kì $0,002\,\text{s}$. Tần số tính theo Hz bằng bao nhiêu? Chọn giá trị đúng.
\choice
{1000}
{250}
{501}
{\True 500}
\loigiai{
Đáp án D. 500. Đối chiếu định nghĩa, hệ thức hoặc dữ kiện của bài toán cho kết luận trên.
}
\end{ex}

% ===== Câu 21 =====
% ID: L11C2B9-06-DS
% Mức: TH
\dangbt{Nhận biết và vận dụng các đặc trưng của sóng âm}
\begin{ex}
% ID: L11C2B9-06-DS
% Mức: TH
Một loa phát âm $500\,\text{Hz}$ trong không khí.
\choiceTF
{\True Tần số âm bằng tần số dao động của nguồn.}
{\True Âm trong không khí là sóng dọc.}
{\True Tăng biên độ nguồn thường làm âm to hơn.}
{Tăng biên độ làm âm cao hơn.}
\loigiai{
\begin{itemchoice}
\item (Đúng) Tần số âm bằng tần số dao động của nguồn. (Vì): Tần số âm bằng tần số dao động của nguồn. b) Đ: Âm trong không khí là sóng dọc. c) Đ: Tăng biên độ nguồn thường làm âm to hơn. d) S: Tăng biên độ làm âm cao hơn
\item (Đúng) Âm trong không khí là sóng dọc.
\item (Đúng) Tăng biên độ nguồn thường làm âm to hơn.
\item (Sai) Tăng biên độ làm âm cao hơn.
\end{itemchoice}
}
\end{ex}

% ===== Câu 22 =====
% ID: L11C2B9-06-TL
% Mức: VD
\dangbt{Bài toán thực tế về âm, độ cao, độ to và môi trường truyền âm}
\begin{ex}
% ID: L11C2B9-06-TL
% Mức: VD
Trong chuyên đề “Bài toán thực tế về âm, độ cao, độ to và môi trường truyền âm”, Mở rộng tình huống sau thành một ví dụ thực tế và trình bày phương pháp giải: Một âm $850\,\text{Hz}$ truyền trong không khí với tốc độ $340\,\text{m}$ /s. Tính bước sóng và cho biết âm thuộc vùng nghe được hay không.
\choiceTF

\loigiai{
$\lambda=340/850=0,4$ m. Tần số $850\,\text{Hz}$ thuộc vùng âm nghe được của người. Có thể kiểm chứng bằng cách thay kết quả vào dữ kiện ban đầu và đối chiếu với hiện tượng thực tế.
}
\end{ex}

% ===== Câu 23 =====
% ID: L11C2B9-06-TLN
% Mức: VD
\begin{ex}
% ID: L11C2B9-06-TLN
% Mức: VD
Trong chuyên đề “Bài toán thực tế về âm, độ cao, độ to và môi trường truyền âm”, Từ kết quả của tình huống sau, nếu đại lượng cần tìm được nhân với hệ số 1.5 thì giá trị mới bằng bao nhiêu? Siêu âm tần số 40 kHz truyền với tốc độ $320\,\text{m}$ /s. Bước sóng tính theo mm bằng bao nhiêu?
\shortans{12}
\loigiai{
Xác định đúng dữ kiện và đơn vị, áp dụng hệ thức của dạng bài rồi tính được kết quả 12.
}
\end{ex}

% ===== Câu 24 =====
% ID: L11C2B9-06-TN
% Mức: TH
\begin{ex}
% ID: L11C2B9-06-TN
% Mức: TH
Trong chuyên đề “Bài toán thực tế về âm, độ cao, độ to và môi trường truyền âm”, Siêu âm tần số 40 kHz truyền với tốc độ $320\,\text{m}$ /s. Bước sóng tính theo mm bằng bao nhiêu? Chọn giá trị đúng.
\choice
{\True 8}
{16}
{4}
{9}
\loigiai{
Đáp án A. 8. Đối chiếu định nghĩa, hệ thức hoặc dữ kiện của bài toán cho kết luận trên.
}
\end{ex}

% ===== Câu 25 =====
% ID: L11C2B9-07-DS
% Mức: TH
\dangbt{Nhận biết và vận dụng các đặc trưng của sóng âm}
\begin{ex}
% ID: L11C2B9-07-DS
% Mức: TH
Xét hạ âm, âm nghe được và siêu âm.
\choiceTF
{\True Hạ âm có tần số dưới ngưỡng nghe thấp.}
{\True Siêu âm có tần số lớn hơn 20 kHz theo quy ước thường dùng.}
{\True Tai người bình thường chỉ nghe được một khoảng tần số hữu hạn.}
{Mọi động vật có cùng giới hạn nghe như người.}
\loigiai{
\begin{itemchoice}
\item (Đúng) Hạ âm có tần số dưới ngưỡng nghe thấp. (Vì): Hạ âm có tần số dưới ngưỡng nghe thấp. b) Đ: Siêu âm có tần số lớn hơn 20 kHz theo quy ước thường dùng. c) Đ: Tai người bình thường chỉ nghe được một khoảng tần số hữu hạn. d) S: Mọi động vật có cùng giới hạn nghe như người
\item (Đúng) Siêu âm có tần số lớn hơn 20 kHz theo quy ước thường dùng.
\item (Đúng) Tai người bình thường chỉ nghe được một khoảng tần số hữu hạn.
\item (Sai) Mọi động vật có cùng giới hạn nghe như người.
\end{itemchoice}
}
\end{ex}

% ===== Câu 26 =====
% ID: L11C2B9-07-TL
% Mức: TH
\dangbt{Chưa có dạng}
\begin{ex}
% ID: L11C2B9-07-TL
% Mức: TH
Một sóng cơ có tần số $20 \mathrm{Hz}$ truyền trên mặt nước với tốc độ $1,5 \mathrm{m/s}$. Trên phương truyền sóng, sóng truyền tới điểm $P$ rồi mới tới điểm $Q$ cách nó $16,125 \mathrm{cm}$. Tại thời điểm t, điểm $P$ hạ xuống thấp nhất thì sau thời gian ngắn nhất là bao nhiêu điểm $Q$ sẽ hạ xuống thấp nhất? $\nguon{ly11 kntt.pdf}$
\choiceTF

\loigiai{
Ta có: $\lambda = \dfrac{v}{f} = \dfrac{1,5}{20}=0,075\,\mathrm{m}=7,5\,\mathrm{cm}PQ =16,125\,\mathrm{cm}$ = 2λ+0,15λ = Q'Q+PQ

[Một sóng cơ có tần số $20\,\mathrm{Hz}$ truyền trên mặt nước với tốc độ 1, $5\,\mathrm{m}$ /s]

Kết hợp với sử dụng đồ thị Hình 9.3 $G$ ta thấy thời gian ngắn nhất để Q' đi từ vị trí hiện tại đến vị trí thấp nhất là 0,15 $T$ = $\dfrac{0,15}{20} = \dfrac{3}{400}\text{s}$ .
}
\end{ex}

% ===== Câu 27 =====
% ID: L11C2B9-07-TLN
% Mức: TH
\dangbt{Nhận biết và vận dụng các đặc trưng của sóng âm}
\begin{ex}
% ID: L11C2B9-07-TLN
% Mức: TH
Âm có tốc độ $340\,\text{m}$ /s và tần số $680\,\text{Hz}$. Bước sóng tính theo m bằng bao nhiêu?
\shortans{0.5}
\loigiai{
Xác định đúng dữ kiện và đơn vị, áp dụng hệ thức của dạng bài rồi tính được kết quả 0.5.
}
\end{ex}

% ===== Câu 28 =====
% ID: L11C2B9-07-TN
% Mức: VD
\dangbt{Bài toán thực tế về âm, độ cao, độ to và môi trường truyền âm}
\begin{ex}
% ID: L11C2B9-07-TN
% Mức: VD
Trong chuyên đề “Bài toán thực tế về âm, độ cao, độ to và môi trường truyền âm”, Kết luận nào phù hợp nhất khi giải bài toán sau: Phân biệt độ cao và độ to của âm bằng hai đặc trưng vật lí tương ứng.
\choice
{Kết quả không phụ thuộc môi trường, tần số hoặc điều kiện biên.}
{\True Độ cao phụ thuộc chủ yếu vào tần số}
{Không cần sử dụng định nghĩa hay dữ kiện đã cho.}
{Đại lượng cần tìm luôn bằng 0 trong mọi trường hợp.}
\loigiai{
Đáp án B. Độ cao phụ thuộc chủ yếu vào tần số. Đối chiếu định nghĩa, hệ thức hoặc dữ kiện của bài toán cho kết luận trên.
}
\end{ex}

% ===== Câu 29 =====
% ID: L11C2B9-08-DS
% Mức: TH
\dangbt{Nhận biết và vận dụng các đặc trưng của sóng âm}
\begin{ex}
% ID: L11C2B9-08-DS
% Mức: TH
Âm truyền từ không khí vào nước.
\choiceTF
{\True Tần số do nguồn quyết định nên không đổi.}
{\True Tốc độ truyền có thể thay đổi.}
{\True Bước sóng thay đổi khi tốc độ thay đổi.}
{Âm truyền được trong chân không giữa hai môi trường.}
\loigiai{
\begin{itemchoice}
\item (Đúng) Tần số do nguồn quyết định nên không đổi. (Vì): Tần số do nguồn quyết định nên không đổi. b) Đ: Tốc độ truyền có thể thay đổi. c) Đ: Bước sóng thay đổi khi tốc độ thay đổi. d) S: Âm truyền được trong chân không giữa hai môi trường
\item (Đúng) Tốc độ truyền có thể thay đổi.
\item (Đúng) Bước sóng thay đổi khi tốc độ thay đổi.
\item (Sai) Âm truyền được trong chân không giữa hai môi trường.
\end{itemchoice}
}
\end{ex}

% ===== Câu 30 =====
% ID: L11C2B9-08-TL
% Mức: TH
\dangbt{Chưa có dạng}
\begin{ex}
% ID: L11C2B9-08-TL
% Mức: TH
*. $P$ và $Q$ là hai điểm trên mặt nước cách nhau một khoảng $20 \mathrm{cm}$. Tại một điểm $O$ trên đường thẳng PQ và nằm ngoài đoạn PQ, người ta đặt nguồn dao động điều hoà theo phương vuông góc với mặt nước với phương trình $u = 5 cos\omega t(\mathrm{cm})$, tạo ra sóng trên mặt nước với bước sóng $\lambda =$ 15 $\mathrm{cm}$. Khoảng cách xa nhất và gần nhất giữa hai phần tử môi trường tại $P$ và $Q$ khi có sóng truyền qua là bao nhiêu? $\nguon{ly11 kntt.pdf}$
\choiceTF

\loigiai{
Đối với trường hợp sóng ngang, khoảng cách giữa hai điểm $P$, $Q$ khi dao động được mô tả như Hình $9.4\,\text{G}$.

[ $P$ và $Q$ là hai điểm trên mặt nước cách nhau một khoảng $20\,\mathrm{cm}$]

Gọi O₁, O₂ lần lượt là vị trí cân bằng của $P$ và $Q$; u₁, u₂ lần lượt là li độ dao động của các phần tử tại $P$ và $Q$; $\Delta$ u = u₁ - u₂.

Khoảng cách giữa $P$ và $Q$ trong quá trình dao động là:

[ $P$ và $Q$ là hai điểm trên mặt nước cách nhau một khoảng $20\,\mathrm{cm}$]

Vậy khoảng cách gần nhất giữa $P$ và $Q$ là: l_(min) = O₁O₂ = $20\,\mathrm{cm}$.

Khoảng cách xa nhất giữa $P$ và $Q$ là: $l_{\text{max}} = \sqrt{\left( {\text{O}_{1}\text{O}_{2}} \right)^{2} + \left( \text{\Delta u}_{\text{max}} \right)^{2}}$.

Giả sử sóng truyền qua $P$ rồi mới đến $Q$ thì dao động tại $P$ sớm pha hơn $Q$ là: $\text{\Delta }\varphi = \dfrac{2\pi\left( {PQ} \right)}{\lambda} = \dfrac{8\pi}{3}$

Chọn mốc thời gian để phương trình dao động của phần tử tại $P$ là: u₁ = 5cosωt(cm)

thì phương trình dao động của phần tử tại $Q$ là: $u_{2} = 5\cos\left( {\omega t - \dfrac{8\pi}{3}} \right)\left( \text{cm} \right)$.

$\text{\Delta u} = \text{u}_{1} - \text{u}_{2} = 5\cos\left( {\omega\text{t} - \dfrac{8\pi}{3}} \right) - 5\cos\omega\text{t} = 5\sqrt{3}\cos\left( {\omega\text{t} - \dfrac{5\pi}{6}} \right)\left( \text{cm} \right)\left. \Rightarrow\text{\Delta u}_{\text{max}} = 5\sqrt{3}\text{cm} \right.$.

$l_{\text{max}} = \sqrt{(20)^{2} + \left( 5\sqrt{3} \right)^{2}} = 5\sqrt{19}\text{cm}.$
}
\end{ex}

% ===== Câu 31 =====
% ID: L11C2B9-08-TLN
% Mức: TH
\dangbt{Nhận biết và vận dụng các đặc trưng của sóng âm}
\begin{ex}
% ID: L11C2B9-08-TLN
% Mức: TH
Một âm có chu kì $0,002\,\text{s}$. Tần số tính theo Hz bằng bao nhiêu?
\shortans{500}
\loigiai{
Xác định đúng dữ kiện và đơn vị, áp dụng hệ thức của dạng bài rồi tính được kết quả 500.
}
\end{ex}

% ===== Câu 32 =====
% ID: L11C2B9-08-TN
% Mức: VD
\dangbt{Bài toán thực tế về âm, độ cao, độ to và môi trường truyền âm}
\begin{ex}
% ID: L11C2B9-08-TN
% Mức: VD
Trong chuyên đề “Bài toán thực tế về âm, độ cao, độ to và môi trường truyền âm”, Kết luận nào phù hợp nhất khi giải bài toán sau: Giải thích vì sao cùng một nốt nhạc truyền từ không khí vào nước vẫn có tần số không đổi nhưng bước sóng thay đổi.
\choice
{Đại lượng cần tìm luôn bằng 0 trong mọi trường hợp.}
{Kết quả không phụ thuộc môi trường, tần số hoặc điều kiện biên.}
{\True Tần số do nguồn quyết định và liên tục qua mặt phân cách}
{Không cần sử dụng định nghĩa hay dữ kiện đã cho.}
\loigiai{
Đáp án C. Tần số do nguồn quyết định và liên tục qua mặt phân cách. Đối chiếu định nghĩa, hệ thức hoặc dữ kiện của bài toán cho kết luận trên.
}
\end{ex}

% ===== Câu 33 =====
% ID: L11C2B9-09-DS
% Mức: VD
\dangbt{Nhận biết và vận dụng các đặc trưng của sóng âm}
\begin{ex}
% ID: L11C2B9-09-DS
% Mức: VD
Xét các nhận định tổng hợp thuộc dạng “Nhận biết và vận dụng các đặc trưng của sóng âm”.
\choiceTF
{\True tần số.}
{\True biên độ dao động âm.}
{\True siêu âm.}
{biên độ.}
\loigiai{
\begin{itemchoice}
\item (Đúng) tần số. (Vì): tần số. b) Đ: biên độ dao động âm. c) Đ: siêu âm. d) S: biên độ
\item (Đúng) biên độ dao động âm.
\item (Đúng) siêu âm.
\item (Sai) biên độ.
\end{itemchoice}
}
\end{ex}

% ===== Câu 34 =====
% ID: L11C2B9-09-TL
% Mức: TH
\dangbt{Chưa có dạng}
\begin{bt}
% ID: L11C2B9-09-TL
% Mức: TH
*. Một sóng dọc truyền trong môi trường với bước sóng $15 \mathrm{cm}$, biên độ không đổi $A =$ 5 $\sqrt{3}\mathrm{cm}$. Gọi $P$ và $Q$ là hai điểm cùng nằm trên một phương truyền sóng. Khi chưa có sóng truyền đến hai điểm $P$ và $Q$ nằm cách nguồn các khoảng lần lượt là $20 \mathrm{cm}$ và $30 \mathrm{cm}$. Khoảng cách xa nhất và gần nhất giữa hai phần tử môi trường tại $P$ và $Q$ khi có sóng truyền qua là bao nhiêu? BÀl 11. SÓNG ĐIỆ $N$ TỪ $\nguon{ly11 kntt.pdf}$
\loigiai{
Đối với trường hợp sóng dọc, khoảng cách giữa hai điểm $P$, $Q$ khi dao động được mô tả như Hình $9.5\,\text{G}$.

[Một sóng dọc truyền trong môi trường với bước sóng $15\,\mathrm{cm}$]

Gọi O₁, O₂ lần lượt là vị trí cân bằng của $P$ và $Q$; u₁, u₂ lần lượt là li độ dao động của các phần tử tại $P$ và $Q$; $\Delta$ u = u₁ - u₂.

Khoảng cách giữa $P$ và $Q$ trong quá trình dao động là:

[Một sóng dọc truyền trong môi trường với bước sóng $15\,\mathrm{cm}$]

Giả sử sóng truyền qua $P$ rồi mới đến $Q$ thì dao động tại $P$ sớm pha hơn $Q$ là: $\text{\Delta }\varphi = \dfrac{2\pi\left( {PQ} \right)}{\lambda} = \dfrac{4\pi}{3}$

Chọn mốc thời gian để phương trình dao của phần tử tại $P$ là: u₁ = 5 $\sqrt{3}$ cosωt(cm)

thì phương trình dao động của phần tử tại $Q$ là: $u_{2} = 5\sqrt{3}\cos\left( {\omega t - \dfrac{4\pi}{3}} \right)\left( \text{cm} \right)$.

$\text{\Delta }u = u_{1} - u_{2} = 5\sqrt{3}\cos\left( {\omega t - \dfrac{4\pi}{3}} \right) - 5\sqrt{3}\cos\omega t = 15\cos\left( {\omega t - \dfrac{5\pi}{6}} \right)\left( \text{cm} \right)$.

$\Rightarrow\Delta$ u_(max)= $15\,\mathrm{cm}$
}
\end{bt}

% ===== Câu 35 =====
% ID: L11C2B9-09-TLN
% Mức: TH
\dangbt{Nhận biết và vận dụng các đặc trưng của sóng âm}
\begin{ex}
% ID: L11C2B9-09-TLN
% Mức: TH
Siêu âm tần số 40 kHz truyền với tốc độ $320\,\text{m}$ /s. Bước sóng tính theo mm bằng bao nhiêu?
\shortans{8}
\loigiai{
Xác định đúng dữ kiện và đơn vị, áp dụng hệ thức của dạng bài rồi tính được kết quả 8.
}
\end{ex}

% ===== Câu 36 =====
% ID: L11C2B9-09-TN
% Mức: VD
\dangbt{Bài toán thực tế về âm, độ cao, độ to và môi trường truyền âm}
\begin{ex}
% ID: L11C2B9-09-TN
% Mức: VD
Trong chuyên đề “Bài toán thực tế về âm, độ cao, độ to và môi trường truyền âm”, Kết luận nào phù hợp nhất khi giải bài toán sau: Một âm $850\,\text{Hz}$ truyền trong không khí với tốc độ $340\,\text{m}$ /s. Tính bước sóng và cho biết âm thuộc vùng nghe được hay không.
\choice
{Không cần sử dụng định nghĩa hay dữ kiện đã cho.}
{Đại lượng cần tìm luôn bằng 0 trong mọi trường hợp.}
{Kết quả không phụ thuộc môi trường, tần số hoặc điều kiện biên.}
{\True $\lambda=340/850=0,4$ m. Tần số $850\,\text{Hz}$ thuộc vùng âm nghe được của người.}
\loigiai{
Đáp án D. $\lambda=340/850=0,4$ m. Tần số $850\,\text{Hz}$ thuộc vùng âm nghe được của người.. Đối chiếu định nghĩa, hệ thức hoặc dữ kiện của bài toán cho kết luận trên.
}
\end{ex}

% ===== Câu 37 =====
% ID: L11C2B9-10-DS
% Mức: VD
\dangbt{Nhận biết và vận dụng các đặc trưng của sóng âm}
\begin{ex}
% ID: L11C2B9-10-DS
% Mức: VD
Đánh giá cách xử lí các dữ kiện định lượng của dạng “Nhận biết và vận dụng các đặc trưng của sóng âm”.
\choiceTF
{\True Kết quả của bài toán thứ nhất là 0.5.}
{\True Kết quả của bài toán thứ hai là 500.}
{\True Kết quả của bài toán thứ ba là 8.}
{Có thể bỏ qua hoàn toàn đơn vị và điều kiện vật lí khi tính toán.}
\loigiai{
\begin{itemchoice}
\item (Đúng) Kết quả của bài toán thứ nhất là 0.5. (Vì): Kết quả của bài toán thứ nhất là 0.5. b) Đ: Kết quả của bài toán thứ hai là 500. c) Đ: Kết quả của bài toán thứ ba là 8. d) S: Có thể bỏ qua hoàn toàn đơn vị và điều kiện vật lí khi tính toán
\item (Đúng) Kết quả của bài toán thứ hai là 500.
\item (Đúng) Kết quả của bài toán thứ ba là 8.
\item (Sai) Có thể bỏ qua hoàn toàn đơn vị và điều kiện vật lí khi tính toán.
\end{itemchoice}
}
\end{ex}

% ===== Câu 38 =====
% ID: L11C2B9-10-TL
% Mức: TH
\begin{ex}
% ID: L11C2B9-10-TL
% Mức: TH
Phân biệt độ cao và độ to của âm bằng hai đặc trưng vật lí tương ứng.
\choiceTF

\loigiai{
Độ cao phụ thuộc chủ yếu vào tần số; độ to liên hệ với biên độ/cường độ âm tại tai.
}
\end{ex}

% ===== Câu 39 =====
% ID: L11C2B9-10-TLN
% Mức: VD
\begin{ex}
% ID: L11C2B9-10-TLN
% Mức: VD
Từ kết quả của tình huống sau, nếu đại lượng cần tìm được nhân với hệ số 2 thì giá trị mới bằng bao nhiêu? Âm có tốc độ $340\,\text{m}$ /s và tần số $680\,\text{Hz}$. Bước sóng tính theo m bằng bao nhiêu?
\shortans{1}
\loigiai{
Xác định đúng dữ kiện và đơn vị, áp dụng hệ thức của dạng bài rồi tính được kết quả 1.
}
\end{ex}

% ===== Câu 40 =====
% ID: L11C2B9-10-TN
% Mức: VD
\dangbt{Bài toán thực tế về âm, độ cao, độ to và môi trường truyền âm}
\begin{ex}
% ID: L11C2B9-10-TN
% Mức: VD
Trong chuyên đề “Bài toán thực tế về âm, độ cao, độ to và môi trường truyền âm”, Phương pháp phù hợp nhất cho dạng “Bài toán thực tế về âm, độ cao, độ to và môi trường truyền âm” là
\choice
{\True xác định dữ kiện, chọn đúng hệ thức hoặc định nghĩa, đổi đơn vị rồi kiểm tra kết quả}
{chỉ thay số mà không cần xét điều kiện áp dụng}
{bỏ qua đơn vị và chiều truyền sóng}
{luôn coi mọi điểm dao động cùng pha}
\loigiai{
Đáp án A. xác định dữ kiện, chọn đúng hệ thức hoặc định nghĩa, đổi đơn vị rồi kiểm tra kết quả. Đối chiếu định nghĩa, hệ thức hoặc dữ kiện của bài toán cho kết luận trên.
}
\end{ex}

% ===== Câu 41 =====
% ID: L11C2B9-11-DS
% Mức: TH
\dangbt{Nhận dạng loại sóng trong các môi trường và thiết bị thực tế}
\begin{ex}
% ID: L11C2B9-11-DS
% Mức: TH
Trong chuyên đề “Nhận dạng loại sóng trong các môi trường và thiết bị thực tế”, So sánh sóng ngang và sóng dọc.
\choiceTF
{\True Cả hai đều có thể là sóng cơ.}
{\True Sóng ngang có phương dao động vuông góc phương truyền.}
{\True Sóng dọc có phương dao động trùng phương truyền.}
{Sóng ngang cơ học truyền được trong mọi chất khí.}
\loigiai{
\begin{itemchoice}
\item (Đúng) Cả hai đều có thể là sóng cơ. (Vì): Cả hai đều có thể là sóng cơ. b) Đ: Sóng ngang có phương dao động vuông góc phương truyền. c) Đ: Sóng dọc có phương dao động trùng phương truyền. d) S: Sóng ngang cơ học truyền được trong mọi chất khí
\item (Đúng) Sóng ngang có phương dao động vuông góc phương truyền.
\item (Đúng) Sóng dọc có phương dao động trùng phương truyền.
\item (Sai) Sóng ngang cơ học truyền được trong mọi chất khí.
\end{itemchoice}
}
\end{ex}

% ===== Câu 42 =====
% ID: L11C2B9-11-TL
% Mức: TH
\dangbt{Nhận biết và vận dụng các đặc trưng của sóng âm}
\begin{ex}
% ID: L11C2B9-11-TL
% Mức: TH
Giải thích vì sao cùng một nốt nhạc truyền từ không khí vào nước vẫn có tần số không đổi nhưng bước sóng thay đổi.
\choiceTF

\loigiai{
Tần số do nguồn quyết định và liên tục qua mặt phân cách; tốc độ phụ thuộc môi trường nên theo $\lambda=v/f$, bước sóng thay đổi.
}
\end{ex}

% ===== Câu 43 =====
% ID: L11C2B9-11-TLN
% Mức: VD
\begin{ex}
% ID: L11C2B9-11-TLN
% Mức: VD
Từ kết quả của tình huống sau, nếu đại lượng cần tìm được nhân với hệ số 0.5 thì giá trị mới bằng bao nhiêu? Một âm có chu kì $0,002\,\text{s}$. Tần số tính theo Hz bằng bao nhiêu?
\shortans{250}
\loigiai{
Xác định đúng dữ kiện và đơn vị, áp dụng hệ thức của dạng bài rồi tính được kết quả 250.
}
\end{ex}

% ===== Câu 44 =====
% ID: L11C2B9-11-TN
% Mức: NB
\dangbt{Chưa có dạng}
\begin{ex}
% ID: L11C2B9-11-TN
% Mức: NB
Chọn câu đúng. $\nguon{ly11 kntt.pdf}$
\choice
{Sóng dọc là sóng truyền dọc theo một sợi dây.}
{Sóng dọc là sóng truyền theo phương thẳng đứng, còn sóng ngang là sóng truyền theo phương nằm ngang.}
{\True Sóng dọc là sóng trong đó phương dao động (của các phần tử môi trường) trùng với phương truyền.}
{Sóng ngang là sóng trong đó phương dao động (của các phần tử môi trường) trùng với phương truyền.}
\loigiai{
Chọn C.

Sóng dọc là sóng trong đó phương dao động (của các phần tử môi trường) trùng với phương truyền.
}
\end{ex}

% ===== Câu 45 =====
% ID: L11C2B9-12-DS
% Mức: TH
\dangbt{Nhận dạng loại sóng trong các môi trường và thiết bị thực tế}
\begin{ex}
% ID: L11C2B9-12-DS
% Mức: TH
Trong chuyên đề “Nhận dạng loại sóng trong các môi trường và thiết bị thực tế”, Một xung nén truyền dọc theo lò xo.
\choiceTF
{\True Đây là sóng dọc.}
{\True Các vòng lò xo dao động dọc theo trục lò xo.}
{\True Vùng nén và dãn lan truyền dọc lò xo.}
{Mỗi vòng lò xo đi từ đầu này tới đầu kia.}
\loigiai{
\begin{itemchoice}
\item (Đúng) Đây là sóng dọc. (Vì): ây là sóng dọc. b) Đ: Các vòng lò xo dao động dọc theo trục lò xo. c) Đ: Vùng nén và dãn lan truyền dọc lò xo. d) S: Mỗi vòng lò xo đi từ đầu này tới đầu kia
\item (Đúng) Các vòng lò xo dao động dọc theo trục lò xo.
\item (Đúng) Vùng nén và dãn lan truyền dọc lò xo.
\item (Sai) Mỗi vòng lò xo đi từ đầu này tới đầu kia.
\end{itemchoice}
}
\end{ex}

% ===== Câu 46 =====
% ID: L11C2B9-12-TL
% Mức: VD
\dangbt{Nhận biết và vận dụng các đặc trưng của sóng âm}
\begin{ex}
% ID: L11C2B9-12-TL
% Mức: VD
Một âm $850\,\text{Hz}$ truyền trong không khí với tốc độ $340\,\text{m}$ /s. Tính bước sóng và cho biết âm thuộc vùng nghe được hay không.
\choiceTF

\loigiai{
$\lambda=340/850=0,4$ m. Tần số $850\,\text{Hz}$ thuộc vùng âm nghe được của người.
}
\end{ex}

% ===== Câu 47 =====
% ID: L11C2B9-12-TLN
% Mức: VD
\begin{ex}
% ID: L11C2B9-12-TLN
% Mức: VD
Từ kết quả của tình huống sau, nếu đại lượng cần tìm được nhân với hệ số 1.5 thì giá trị mới bằng bao nhiêu? Siêu âm tần số 40 kHz truyền với tốc độ $320\,\text{m}$ /s. Bước sóng tính theo mm bằng bao nhiêu?
\shortans{12}
\loigiai{
Xác định đúng dữ kiện và đơn vị, áp dụng hệ thức của dạng bài rồi tính được kết quả 12.
}
\end{ex}

% ===== Câu 48 =====
% ID: L11C2B9-12-TN
% Mức: NB
\dangbt{Chưa có dạng}
\begin{ex}
% ID: L11C2B9-12-TN
% Mức: NB
Sóng cơ không truyền được trong $\nguon{ly11 kntt.pdf}$
\choice
{\True chân không.}
{không khí.}
{nước.}
{kim loại.}
\loigiai{
Chọn A.

Sóng cơ không truyền được trong chân không.
}
\end{ex}

% ===== Câu 49 =====
% ID: L11C2B9-13-DS
% Mức: TH
\dangbt{Nhận dạng loại sóng trong các môi trường và thiết bị thực tế}
\begin{ex}
% ID: L11C2B9-13-DS
% Mức: TH
Trong chuyên đề “Nhận dạng loại sóng trong các môi trường và thiết bị thực tế”, Sóng truyền trên mặt nước trong thí nghiệm lớp học.
\choiceTF
{\True Có thể coi gần đúng là sóng ngang.}
{\True Phần tử mặt nước dao động chủ yếu lên xuống.}
{\True Phương truyền sóng nằm dọc mặt nước.}
{Sóng này truyền được trong chân không.}
\loigiai{
\begin{itemchoice}
\item (Đúng) Có thể coi gần đúng là sóng ngang. (Vì): Có thể coi gần đúng là sóng ngang. b) Đ: Phần tử mặt nước dao động chủ yếu lên xuống. c) Đ: Phương truyền sóng nằm dọc mặt nước. d) S: Sóng này truyền được trong chân không
\item (Đúng) Phần tử mặt nước dao động chủ yếu lên xuống.
\item (Đúng) Phương truyền sóng nằm dọc mặt nước.
\item (Sai) Sóng này truyền được trong chân không.
\end{itemchoice}
}
\end{ex}

% ===== Câu 50 =====
% ID: L11C2B9-13-TL
% Mức: VD
\dangbt{Nhận biết và vận dụng các đặc trưng của sóng âm}
\begin{ex}
% ID: L11C2B9-13-TL
% Mức: VD
Phân tích các bước lập luận và chỉ ra điều kiện áp dụng trong bài toán: Phân biệt độ cao và độ to của âm bằng hai đặc trưng vật lí tương ứng.
\choiceTF

\loigiai{
Độ cao phụ thuộc chủ yếu vào tần số; độ to liên hệ với biên độ/cường độ âm tại tai. Cần nêu rõ điều kiện áp dụng, đổi đơn vị thống nhất và kiểm tra ý nghĩa vật lí của kết quả.
}
\end{ex}

% ===== Câu 51 =====
% ID: L11C2B9-13-TLN
% Mức: TH
\dangbt{Nhận dạng loại sóng trong các môi trường và thiết bị thực tế}
\begin{ex}
% ID: L11C2B9-13-TLN
% Mức: TH
Trong chuyên đề “Nhận dạng loại sóng trong các môi trường và thiết bị thực tế”, Một sóng âm truyền trong không khí với tốc độ $340\,\text{m}$ /s trong $0,5\,\text{s}$. Quãng đường tính theo m bằng bao nhiêu?
\shortans{170}
\loigiai{
Xác định đúng dữ kiện và đơn vị, áp dụng hệ thức của dạng bài rồi tính được kết quả 170.
}
\end{ex}

% ===== Câu 52 =====
% ID: L11C2B9-13-TN
% Mức: NB
\dangbt{Chưa có dạng}
\begin{ex}
% ID: L11C2B9-13-TN
% Mức: NB
Một sóng ngang có tần số $100 \mathrm{Hz}$ truyền trên một sợi dây nằm ngang với tốc độ $60 \mathrm{m/s}$, qua điểm $A$ rồi đến điểm $B$ cách nhau $7,95\,\text{m}$. Tại một thời điểm nào đó $A$ có li độ âm và đang chuyển động đi lên thì điểm $B$ đang có li độ $\nguon{ly11 kntt.pdf}$
\choice
{\True âm và đang đi xuống.}
{âm và đang đi lên.}
{dương và đang đi lên.}
{dương và đang đi xuống}
\loigiai{
Chọn A.

Sử dụng đồ thị li độ quãng đường Hình 9.1 $G$ của sóng quy ước chiều truyền dương để xác định các vùng mà các phần tử vật chất đang đi lên và đi xuống.

[Một sóng ngang có tần số $10\,\mathrm{Hz}$ truyền trên một sợi dây]

Ta có: $\lambda = \dfrac{v}{f} = \dfrac{60}{100}$ = $0,6\,\mathrm{m}$; AB = $7,95\,\mathrm{m}$ = 7,8+0,15 = 13.0,6+0,15 = 13λ+ $\dfrac{\lambda}{4}$.

Từ Hình $9.1\,\text{G}$. Ta thấy $B$ có li độ âm và đang đi xuống.
}
\end{ex}

% ===== Câu 53 =====
% ID: L11C2B9-14-DS
% Mức: VD
\dangbt{Nhận dạng loại sóng trong các môi trường và thiết bị thực tế}
\begin{ex}
% ID: L11C2B9-14-DS
% Mức: VD
Trong chuyên đề “Nhận dạng loại sóng trong các môi trường và thiết bị thực tế”, Xét các nhận định tổng hợp thuộc dạng “Nhận dạng loại sóng trong các môi trường và thiết bị thực tế”.
\choiceTF
{\True sóng ngang.}
{\True sóng dọc.}
{\True chất rắn và trên bề mặt chất lỏng.}
{sóng dọc.}
\loigiai{
\begin{itemchoice}
\item (Đúng) sóng ngang. (Vì): sóng ngang. b) Đ: sóng dọc. c) Đ: chất rắn và trên bề mặt chất lỏng. d) S: sóng dọc
\item (Đúng) sóng dọc.
\item (Đúng) chất rắn và trên bề mặt chất lỏng.
\item (Sai) sóng dọc.
\end{itemchoice}
}
\end{ex}

% ===== Câu 54 =====
% ID: L11C2B9-14-TL
% Mức: VD
\dangbt{Nhận biết và vận dụng các đặc trưng của sóng âm}
\begin{ex}
% ID: L11C2B9-14-TL
% Mức: VD
Một học sinh giải bài sau nhưng bỏ qua điều kiện vật lí. Hãy sửa cách làm và trình bày lời giải đúng: Giải thích vì sao cùng một nốt nhạc truyền từ không khí vào nước vẫn có tần số không đổi nhưng bước sóng thay đổi.
\choiceTF

\loigiai{
Cách giải đúng: Tần số do nguồn quyết định và liên tục qua mặt phân cách; tốc độ phụ thuộc môi trường nên theo $\lambda=v/f$, bước sóng thay đổi. Sai lầm cần tránh là dùng hệ thức khi chưa xác định điều kiện biên, môi trường hoặc đơn vị.
}
\end{ex}

% ===== Câu 55 =====
% ID: L11C2B9-14-TLN
% Mức: TH
\dangbt{Nhận dạng loại sóng trong các môi trường và thiết bị thực tế}
\begin{ex}
% ID: L11C2B9-14-TLN
% Mức: TH
Trong chuyên đề “Nhận dạng loại sóng trong các môi trường và thiết bị thực tế”, Sóng dọc có tần số $200\,\text{Hz}$ và tốc độ $300\,\text{m}$ /s. Bước sóng tính theo m bằng bao nhiêu?
\shortans{01/05/2026}
\loigiai{
Xác định đúng dữ kiện và đơn vị, áp dụng hệ thức của dạng bài rồi tính được kết quả 1.5.
}
\end{ex}

% ===== Câu 56 =====
% ID: L11C2B9-14-TN
% Mức: NB
\dangbt{Chưa có dạng}
\begin{ex}
% ID: L11C2B9-14-TN
% Mức: NB
Một sóng ngang truyền trên một sợi dây rất dài từ $P$ đến Q. Hai điểm $P$, $Q$ trên phương truyền 5 $\lambda$ sóng cách nhau $P$ Q = 4 $. Kết luận nào sau đây là đúng?$ \nguon{ly11kntt.pdf}
\choice
{\True Khi $P$ có li độ cực đại thì $Q$ có vận tốc cực đại.}
{Li độ $P$, $Q$ luôn trái dấu.}
{Khi $Q$ có li độ cực đại thì $P$ có vận tốc cực đại.}
{Khi $P$ có li độ cực đại thì $Q$ qua vị trí cân bằng theo chiều âm. Khi $Q$ có li độ cực đại thì $P$ qua vị trí cân bằng theo chiều dương.}
\loigiai{
Chọn A.

Sử dụng đồ thị li độ - quãng đường Hình 9.2Ga, b.

[Một sóng ngang truyền trên một sợi dây rất dài từ $P$ đến Q]

Ta thấy theo Hình 9.2Ga, khi $Q$ có li độ cực đại thì $P$ qua vị trí cân bằng theo chiều âm; theo Hình 9.2Gb thì khi $P$ có li độ cực đại thì $Q$ qua vị trí cân bằng theo chiều dương.
}
\end{ex}

% ===== Câu 57 =====
% ID: L11C2B9-15-DS
% Mức: VD
\dangbt{Nhận dạng loại sóng trong các môi trường và thiết bị thực tế}
\begin{ex}
% ID: L11C2B9-15-DS
% Mức: VD
Trong chuyên đề “Nhận dạng loại sóng trong các môi trường và thiết bị thực tế”, Đánh giá cách xử lí các dữ kiện định lượng của dạng “Nhận dạng loại sóng trong các môi trường và thiết bị thực tế”.
\choiceTF
{\True Kết quả của bài toán thứ nhất là 170.}
{\True Kết quả của bài toán thứ hai là 1.5.}
{\True Kết quả của bài toán thứ ba là 4.}
{Có thể bỏ qua hoàn toàn đơn vị và điều kiện vật lí khi tính toán.}
\loigiai{
\begin{itemchoice}
\item (Đúng) Kết quả của bài toán thứ nhất là 170. (Vì): Kết quả của bài toán thứ nhất là 170. b) Đ: Kết quả của bài toán thứ hai là 1.5. c) Đ: Kết quả của bài toán thứ ba là 4. d) S: Có thể bỏ qua hoàn toàn đơn vị và điều kiện vật lí khi tính toán
\item (Đúng) Kết quả của bài toán thứ hai là 1.5.
\item (Đúng) Kết quả của bài toán thứ ba là 4.
\item (Sai) Có thể bỏ qua hoàn toàn đơn vị và điều kiện vật lí khi tính toán.
\end{itemchoice}
}
\end{ex}

% ===== Câu 58 =====
% ID: L11C2B9-15-TL
% Mức: VD
\dangbt{Nhận biết và vận dụng các đặc trưng của sóng âm}
\begin{ex}
% ID: L11C2B9-15-TL
% Mức: VD
Mở rộng tình huống sau thành một ví dụ thực tế và trình bày phương pháp giải: Một âm $850\,\text{Hz}$ truyền trong không khí với tốc độ $340\,\text{m}$ /s. Tính bước sóng và cho biết âm thuộc vùng nghe được hay không.
\choiceTF

\loigiai{
$\lambda=340/850=0,4$ m. Tần số $850\,\text{Hz}$ thuộc vùng âm nghe được của người. Có thể kiểm chứng bằng cách thay kết quả vào dữ kiện ban đầu và đối chiếu với hiện tượng thực tế.
}
\end{ex}

% ===== Câu 59 =====
% ID: L11C2B9-15-TLN
% Mức: TH
\dangbt{Nhận dạng loại sóng trong các môi trường và thiết bị thực tế}
\begin{ex}
% ID: L11C2B9-15-TLN
% Mức: TH
Trong chuyên đề “Nhận dạng loại sóng trong các môi trường và thiết bị thực tế”, Sóng trên dây truyền $12\,\text{m}$ trong $3\,\text{s}$. Tốc độ tính theo m/s bằng bao nhiêu?
\shortans{4}
\loigiai{
Xác định đúng dữ kiện và đơn vị, áp dụng hệ thức của dạng bài rồi tính được kết quả 4.
}
\end{ex}

% ===== Câu 60 =====
% ID: L11C2B9-15-TN
% Mức: TH
\dangbt{Chưa có dạng}
\begin{ex}
% ID: L11C2B9-15-TN
% Mức: TH
Tìm phát biểu sai khi nói về sóng cơ. $\nguon{ly11 kntt.pdf}$
\choice
{\True Bước sóng là khoảng cách giữa hai điểm gần nhau nhất trên cùng một phương truyền sóng dao động ngược pha nhau.}
{Sóng trong đó các phần tử môi trường dao động theo phương trùng với phương truyền sóng được gọi là sóng dọc.}
{Tại mỗi điểm của môi trường có sóng truyền qua, biên độ của sóng là biên độ dao động của phần tử môi trường.}
{Sóng trong đó các phần tử môi trường dao động theo phương vuông góc với phương truyền sóng được gọi là sóng ngang.}
\loigiai{
Chọn A.

Bước sóng là khoảng cách giữa hai điểm gần nhau nhất trên cùng một phương truyền sóng dao động cùng pha nhau.
}
\end{ex}

% ===== Câu 61 =====
% ID: L11C2B9-16-DS
% Mức: TH
\dangbt{Phân biệt sóng ngang, sóng dọc và môi trường truyền sóng}
\begin{ex}
% ID: L11C2B9-16-DS
% Mức: TH
So sánh sóng ngang và sóng dọc.
\choiceTF
{\True Cả hai đều có thể là sóng cơ.}
{\True Sóng ngang có phương dao động vuông góc phương truyền.}
{\True Sóng dọc có phương dao động trùng phương truyền.}
{Sóng ngang cơ học truyền được trong mọi chất khí.}
\loigiai{
\begin{itemchoice}
\item (Đúng) Cả hai đều có thể là sóng cơ. (Vì): Cả hai đều có thể là sóng cơ. b) Đ: Sóng ngang có phương dao động vuông góc phương truyền. c) Đ: Sóng dọc có phương dao động trùng phương truyền. d) S: Sóng ngang cơ học truyền được trong mọi chất khí
\item (Đúng) Sóng ngang có phương dao động vuông góc phương truyền.
\item (Đúng) Sóng dọc có phương dao động trùng phương truyền.
\item (Sai) Sóng ngang cơ học truyền được trong mọi chất khí.
\end{itemchoice}
}
\end{ex}

% ===== Câu 62 =====
% ID: L11C2B9-16-TL
% Mức: TH
\dangbt{Nhận dạng loại sóng trong các môi trường và thiết bị thực tế}
\begin{ex}
% ID: L11C2B9-16-TL
% Mức: TH
Trong chuyên đề “Nhận dạng loại sóng trong các môi trường và thiết bị thực tế”, So sánh sóng ngang và sóng dọc theo phương dao động và môi trường truyền.
\choiceTF

\loigiai{
Sóng ngang: phương dao động vuông góc phương truyền, truyền trong chất rắn và trên bề mặt chất lỏng. Sóng dọc: phương dao động trùng phương truyền, truyền trong rắn, lỏng và khí.
}
\end{ex}

% ===== Câu 63 =====
% ID: L11C2B9-16-TLN
% Mức: VD
\begin{ex}
% ID: L11C2B9-16-TLN
% Mức: VD
Trong chuyên đề “Nhận dạng loại sóng trong các môi trường và thiết bị thực tế”, Từ kết quả của tình huống sau, nếu đại lượng cần tìm được nhân với hệ số 2 thì giá trị mới bằng bao nhiêu? Một sóng âm truyền trong không khí với tốc độ $340\,\text{m}$ /s trong $0,5\,\text{s}$. Quãng đường tính theo m bằng bao nhiêu?
\shortans{340}
\loigiai{
Xác định đúng dữ kiện và đơn vị, áp dụng hệ thức của dạng bài rồi tính được kết quả 340.
}
\end{ex}

% ===== Câu 64 =====
% ID: L11C2B9-16-TN
% Mức: NB
\dangbt{Nhận biết và vận dụng các đặc trưng của sóng âm}
\begin{ex}
% ID: L11C2B9-16-TN
% Mức: NB
Độ cao của âm chủ yếu phụ thuộc vào
\choice
{\True tần số}
{biên độ}
{tốc độ gió}
{khoảng cách tới nguồn}
\loigiai{
Đáp án A. tần số. Đối chiếu định nghĩa, hệ thức hoặc dữ kiện của bài toán cho kết luận trên.
}
\end{ex}

% ===== Câu 65 =====
% ID: L11C2B9-17-DS
% Mức: TH
\dangbt{Phân biệt sóng ngang, sóng dọc và môi trường truyền sóng}
\begin{ex}
% ID: L11C2B9-17-DS
% Mức: TH
Một xung nén truyền dọc theo lò xo.
\choiceTF
{\True Đây là sóng dọc.}
{\True Các vòng lò xo dao động dọc theo trục lò xo.}
{\True Vùng nén và dãn lan truyền dọc lò xo.}
{Mỗi vòng lò xo đi từ đầu này tới đầu kia.}
\loigiai{
\begin{itemchoice}
\item (Đúng) Đây là sóng dọc. (Vì): ây là sóng dọc. b) Đ: Các vòng lò xo dao động dọc theo trục lò xo. c) Đ: Vùng nén và dãn lan truyền dọc lò xo. d) S: Mỗi vòng lò xo đi từ đầu này tới đầu kia
\item (Đúng) Các vòng lò xo dao động dọc theo trục lò xo.
\item (Đúng) Vùng nén và dãn lan truyền dọc lò xo.
\item (Sai) Mỗi vòng lò xo đi từ đầu này tới đầu kia.
\end{itemchoice}
}
\end{ex}

% ===== Câu 66 =====
% ID: L11C2B9-17-TL
% Mức: TH
\dangbt{Nhận dạng loại sóng trong các môi trường và thiết bị thực tế}
\begin{ex}
% ID: L11C2B9-17-TL
% Mức: TH
Trong chuyên đề “Nhận dạng loại sóng trong các môi trường và thiết bị thực tế”, Giải thích vì sao sóng âm trong không khí là sóng dọc.
\choiceTF

\loigiai{
Các lớp không khí dao động nén-dãn theo đúng phương truyền âm, tạo các vùng áp suất cao và thấp lan truyền.
}
\end{ex}

% ===== Câu 67 =====
% ID: L11C2B9-17-TLN
% Mức: VD
\begin{ex}
% ID: L11C2B9-17-TLN
% Mức: VD
Trong chuyên đề “Nhận dạng loại sóng trong các môi trường và thiết bị thực tế”, Từ kết quả của tình huống sau, nếu đại lượng cần tìm được nhân với hệ số 0.5 thì giá trị mới bằng bao nhiêu? Sóng dọc có tần số $200\,\text{Hz}$ và tốc độ $300\,\text{m}$ /s. Bước sóng tính theo m bằng bao nhiêu?
\shortans{0.75}
\loigiai{
Xác định đúng dữ kiện và đơn vị, áp dụng hệ thức của dạng bài rồi tính được kết quả 0.75.
}
\end{ex}

% ===== Câu 68 =====
% ID: L11C2B9-17-TN
% Mức: NB
\dangbt{Nhận biết và vận dụng các đặc trưng của sóng âm}
\begin{ex}
% ID: L11C2B9-17-TN
% Mức: NB
Độ to của âm liên hệ trực tiếp nhất với
\choice
{pha ban đầu}
{\True biên độ dao động âm}
{chu kì âm}
{bước sóng duy nhất}
\loigiai{
Đáp án B. biên độ dao động âm. Đối chiếu định nghĩa, hệ thức hoặc dữ kiện của bài toán cho kết luận trên.
}
\end{ex}

% ===== Câu 69 =====
% ID: L11C2B9-18-DS
% Mức: TH
\dangbt{Phân biệt sóng ngang, sóng dọc và môi trường truyền sóng}
\begin{ex}
% ID: L11C2B9-18-DS
% Mức: TH
Sóng truyền trên mặt nước trong thí nghiệm lớp học.
\choiceTF
{\True Có thể coi gần đúng là sóng ngang.}
{\True Phần tử mặt nước dao động chủ yếu lên xuống.}
{\True Phương truyền sóng nằm dọc mặt nước.}
{Sóng này truyền được trong chân không.}
\loigiai{
\begin{itemchoice}
\item (Đúng) Có thể coi gần đúng là sóng ngang. (Vì): Có thể coi gần đúng là sóng ngang. b) Đ: Phần tử mặt nước dao động chủ yếu lên xuống. c) Đ: Phương truyền sóng nằm dọc mặt nước. d) S: Sóng này truyền được trong chân không
\item (Đúng) Phần tử mặt nước dao động chủ yếu lên xuống.
\item (Đúng) Phương truyền sóng nằm dọc mặt nước.
\item (Sai) Sóng này truyền được trong chân không.
\end{itemchoice}
}
\end{ex}

% ===== Câu 70 =====
% ID: L11C2B9-18-TL
% Mức: VD
\dangbt{Nhận dạng loại sóng trong các môi trường và thiết bị thực tế}
\begin{ex}
% ID: L11C2B9-18-TL
% Mức: VD
Trong chuyên đề “Nhận dạng loại sóng trong các môi trường và thiết bị thực tế”, Nêu cách dùng hướng chuyển động của một phần tử để nhận biết loại sóng trên dây hoặc lò xo.
\choiceTF

\loigiai{
So sánh hướng dao động phần tử với hướng truyền: vuông góc là sóng ngang, song song/trùng phương là sóng dọc.
}
\end{ex}

% ===== Câu 71 =====
% ID: L11C2B9-18-TLN
% Mức: VD
\begin{ex}
% ID: L11C2B9-18-TLN
% Mức: VD
Trong chuyên đề “Nhận dạng loại sóng trong các môi trường và thiết bị thực tế”, Từ kết quả của tình huống sau, nếu đại lượng cần tìm được nhân với hệ số 1.5 thì giá trị mới bằng bao nhiêu? Sóng trên dây truyền $12\,\text{m}$ trong $3\,\text{s}$. Tốc độ tính theo m/s bằng bao nhiêu?
\shortans{6}
\loigiai{
Xác định đúng dữ kiện và đơn vị, áp dụng hệ thức của dạng bài rồi tính được kết quả 6.
}
\end{ex}

% ===== Câu 72 =====
% ID: L11C2B9-18-TN
% Mức: TH
\dangbt{Nhận biết và vận dụng các đặc trưng của sóng âm}
\begin{ex}
% ID: L11C2B9-18-TN
% Mức: TH
Âm có tần số 25 kHz được gọi là
\choice
{âm nghe được}
{sóng vô tuyến}
{\True siêu âm}
{hạ âm}
\loigiai{
Đáp án C. siêu âm. Đối chiếu định nghĩa, hệ thức hoặc dữ kiện của bài toán cho kết luận trên.
}
\end{ex}

% ===== Câu 73 =====
% ID: L11C2B9-19-DS
% Mức: VD
\dangbt{Phân biệt sóng ngang, sóng dọc và môi trường truyền sóng}
\begin{ex}
% ID: L11C2B9-19-DS
% Mức: VD
Xét các nhận định tổng hợp thuộc dạng “Phân biệt sóng ngang, sóng dọc và môi trường truyền sóng”.
\choiceTF
{\True sóng ngang.}
{\True sóng dọc.}
{\True chất rắn và trên bề mặt chất lỏng.}
{sóng dọc.}
\loigiai{
\begin{itemchoice}
\item (Đúng) sóng ngang. (Vì): sóng ngang. b) Đ: sóng dọc. c) Đ: chất rắn và trên bề mặt chất lỏng. d) S: sóng dọc
\item (Đúng) sóng dọc.
\item (Đúng) chất rắn và trên bề mặt chất lỏng.
\item (Sai) sóng dọc.
\end{itemchoice}
}
\end{ex}

% ===== Câu 74 =====
% ID: L11C2B9-19-TL
% Mức: VD
\dangbt{Nhận dạng loại sóng trong các môi trường và thiết bị thực tế}
\begin{ex}
% ID: L11C2B9-19-TL
% Mức: VD
Trong chuyên đề “Nhận dạng loại sóng trong các môi trường và thiết bị thực tế”, Phân tích các bước lập luận và chỉ ra điều kiện áp dụng trong bài toán: So sánh sóng ngang và sóng dọc theo phương dao động và môi trường truyền.
\choiceTF

\loigiai{
Sóng ngang: phương dao động vuông góc phương truyền, truyền trong chất rắn và trên bề mặt chất lỏng. Sóng dọc: phương dao động trùng phương truyền, truyền trong rắn, lỏng và khí. Cần nêu rõ điều kiện áp dụng, đổi đơn vị thống nhất và kiểm tra ý nghĩa vật lí của kết quả.
}
\end{ex}

% ===== Câu 75 =====
% ID: L11C2B9-19-TLN
% Mức: TH
\dangbt{Phân biệt sóng ngang, sóng dọc và môi trường truyền sóng}
\begin{ex}
% ID: L11C2B9-19-TLN
% Mức: TH
Một sóng âm truyền trong không khí với tốc độ $340\,\text{m}$ /s trong $0,5\,\text{s}$. Quãng đường tính theo m bằng bao nhiêu?
\shortans{170}
\loigiai{
Xác định đúng dữ kiện và đơn vị, áp dụng hệ thức của dạng bài rồi tính được kết quả 170.
}
\end{ex}

% ===== Câu 76 =====
% ID: L11C2B9-19-TN
% Mức: TH
\dangbt{Nhận biết và vận dụng các đặc trưng của sóng âm}
\begin{ex}
% ID: L11C2B9-19-TN
% Mức: TH
Âm có tốc độ $340\,\text{m}$ /s và tần số $680\,\text{Hz}$. Bước sóng tính theo m bằng bao nhiêu? Chọn giá trị đúng.
\choice
{1}
{0.25}
{01/05/2026}
{\True 0.5}
\loigiai{
Đáp án D. 0.5. Đối chiếu định nghĩa, hệ thức hoặc dữ kiện của bài toán cho kết luận trên.
}
\end{ex}

% ===== Câu 77 =====
% ID: L11C2B9-20-DS
% Mức: VD
\dangbt{Phân biệt sóng ngang, sóng dọc và môi trường truyền sóng}
\begin{ex}
% ID: L11C2B9-20-DS
% Mức: VD
Đánh giá cách xử lí các dữ kiện định lượng của dạng “Phân biệt sóng ngang, sóng dọc và môi trường truyền sóng”.
\choiceTF
{\True Kết quả của bài toán thứ nhất là 170.}
{\True Kết quả của bài toán thứ hai là 1.5.}
{\True Kết quả của bài toán thứ ba là 4.}
{Có thể bỏ qua hoàn toàn đơn vị và điều kiện vật lí khi tính toán.}
\loigiai{
\begin{itemchoice}
\item (Đúng) Kết quả của bài toán thứ nhất là 170. (Vì): Kết quả của bài toán thứ nhất là 170. b) Đ: Kết quả của bài toán thứ hai là 1.5. c) Đ: Kết quả của bài toán thứ ba là 4. d) S: Có thể bỏ qua hoàn toàn đơn vị và điều kiện vật lí khi tính toán
\item (Đúng) Kết quả của bài toán thứ hai là 1.5.
\item (Đúng) Kết quả của bài toán thứ ba là 4.
\item (Sai) Có thể bỏ qua hoàn toàn đơn vị và điều kiện vật lí khi tính toán.
\end{itemchoice}
}
\end{ex}

% ===== Câu 78 =====
% ID: L11C2B9-20-TL
% Mức: VD
\dangbt{Nhận dạng loại sóng trong các môi trường và thiết bị thực tế}
\begin{ex}
% ID: L11C2B9-20-TL
% Mức: VD
Trong chuyên đề “Nhận dạng loại sóng trong các môi trường và thiết bị thực tế”, Một học sinh giải bài sau nhưng bỏ qua điều kiện vật lí. Hãy sửa cách làm và trình bày lời giải đúng: Giải thích vì sao sóng âm trong không khí là sóng dọc.
\choiceTF

\loigiai{
Cách giải đúng: Các lớp không khí dao động nén-dãn theo đúng phương truyền âm, tạo các vùng áp suất cao và thấp lan truyền. Sai lầm cần tránh là dùng hệ thức khi chưa xác định điều kiện biên, môi trường hoặc đơn vị.
}
\end{ex}

% ===== Câu 79 =====
% ID: L11C2B9-20-TLN
% Mức: TH
\dangbt{Phân biệt sóng ngang, sóng dọc và môi trường truyền sóng}
\begin{ex}
% ID: L11C2B9-20-TLN
% Mức: TH
Sóng dọc có tần số $200\,\text{Hz}$ và tốc độ $300\,\text{m}$ /s. Bước sóng tính theo m bằng bao nhiêu?
\shortans{01/05/2026}
\loigiai{
Xác định đúng dữ kiện và đơn vị, áp dụng hệ thức của dạng bài rồi tính được kết quả 1.5.
}
\end{ex}

% ===== Câu 80 =====
% ID: L11C2B9-20-TN
% Mức: TH
\dangbt{Nhận biết và vận dụng các đặc trưng của sóng âm}
\begin{ex}
% ID: L11C2B9-20-TN
% Mức: TH
Một âm có chu kì $0,002\,\text{s}$. Tần số tính theo Hz bằng bao nhiêu? Chọn giá trị đúng.
\choice
{\True 500}
{1000}
{250}
{501}
\loigiai{
Đáp án A. 500. Đối chiếu định nghĩa, hệ thức hoặc dữ kiện của bài toán cho kết luận trên.
}
\end{ex}

% ===== Câu 81 =====
% ID: L11C2B9-21-DS
% Mức: TH
\dangbt{Phân tích sự truyền năng lượng của sóng cơ}
\begin{ex}
% ID: L11C2B9-21-DS
% Mức: TH
Một nguồn rung tạo sóng trên dây.
\choiceTF
{\True Nguồn thực hiện công lên dây.}
{\True Năng lượng truyền dọc theo dây.}
{\True Các phần tử dây nhận năng lượng và dao động.}
{Khối lượng dây được chuyển liên tục từ nguồn tới đầu kia.}
\loigiai{
\begin{itemchoice}
\item (Đúng) Nguồn thực hiện công lên dây. (Vì): Nguồn thực hiện công lên dây. b) Đ: Năng lượng truyền dọc theo dây. c) Đ: Các phần tử dây nhận năng lượng và dao động. d) S: Khối lượng dây được chuyển liên tục từ nguồn tới đầu kia
\item (Đúng) Năng lượng truyền dọc theo dây.
\item (Đúng) Các phần tử dây nhận năng lượng và dao động.
\item (Sai) Khối lượng dây được chuyển liên tục từ nguồn tới đầu kia.
\end{itemchoice}
}
\end{ex}

% ===== Câu 82 =====
% ID: L11C2B9-21-TL
% Mức: VD
\dangbt{Nhận dạng loại sóng trong các môi trường và thiết bị thực tế}
\begin{ex}
% ID: L11C2B9-21-TL
% Mức: VD
Trong chuyên đề “Nhận dạng loại sóng trong các môi trường và thiết bị thực tế”, Mở rộng tình huống sau thành một ví dụ thực tế và trình bày phương pháp giải: Nêu cách dùng hướng chuyển động của một phần tử để nhận biết loại sóng trên dây hoặc lò xo.
\choiceTF

\loigiai{
So sánh hướng dao động phần tử với hướng truyền: vuông góc là sóng ngang, song song/trùng phương là sóng dọc. Có thể kiểm chứng bằng cách thay kết quả vào dữ kiện ban đầu và đối chiếu với hiện tượng thực tế.
}
\end{ex}

% ===== Câu 83 =====
% ID: L11C2B9-21-TLN
% Mức: TH
\dangbt{Phân biệt sóng ngang, sóng dọc và môi trường truyền sóng}
\begin{ex}
% ID: L11C2B9-21-TLN
% Mức: TH
Sóng trên dây truyền $12\,\text{m}$ trong $3\,\text{s}$. Tốc độ tính theo m/s bằng bao nhiêu?
\shortans{4}
\loigiai{
Xác định đúng dữ kiện và đơn vị, áp dụng hệ thức của dạng bài rồi tính được kết quả 4.
}
\end{ex}

% ===== Câu 84 =====
% ID: L11C2B9-21-TN
% Mức: TH
\dangbt{Nhận biết và vận dụng các đặc trưng của sóng âm}
\begin{ex}
% ID: L11C2B9-21-TN
% Mức: TH
Siêu âm tần số 40 kHz truyền với tốc độ $320\,\text{m}$ /s. Bước sóng tính theo mm bằng bao nhiêu? Chọn giá trị đúng.
\choice
{9}
{\True 8}
{16}
{4}
\loigiai{
Đáp án B. 8. Đối chiếu định nghĩa, hệ thức hoặc dữ kiện của bài toán cho kết luận trên.
}
\end{ex}

% ===== Câu 85 =====
% ID: L11C2B9-22-DS
% Mức: TH
\dangbt{Phân tích sự truyền năng lượng của sóng cơ}
\begin{ex}
% ID: L11C2B9-22-DS
% Mức: TH
So sánh hai sóng cùng môi trường, cùng tần số nhưng biên độ khác nhau.
\choiceTF
{\True Sóng có biên độ lớn hơn thường mang nhiều năng lượng hơn.}
{\True Cả hai có cùng tốc độ truyền nếu môi trường không đổi.}
{\True Tần số do nguồn quyết định.}
{Biên độ không liên quan mức độ dao động của phần tử.}
\loigiai{
\begin{itemchoice}
\item (Đúng) Sóng có biên độ lớn hơn thường mang nhiều năng lượng hơn. (Vì): Sóng có biên độ lớn hơn thường mang nhiều năng lượng hơn. b) Đ: Cả hai có cùng tốc độ truyền nếu môi trường không đổi. c) Đ: Tần số do nguồn quyết định. d) S: Biên độ không liên quan mức độ dao động của phần tử
\item (Đúng) Cả hai có cùng tốc độ truyền nếu môi trường không đổi.
\item (Đúng) Tần số do nguồn quyết định.
\item (Sai) Biên độ không liên quan mức độ dao động của phần tử.
\end{itemchoice}
}
\end{ex}

% ===== Câu 86 =====
% ID: L11C2B9-22-TL
% Mức: TH
\dangbt{Phân biệt sóng ngang, sóng dọc và môi trường truyền sóng}
\begin{ex}
% ID: L11C2B9-22-TL
% Mức: TH
So sánh sóng ngang và sóng dọc theo phương dao động và môi trường truyền.
\choiceTF

\loigiai{
Sóng ngang: phương dao động vuông góc phương truyền, truyền trong chất rắn và trên bề mặt chất lỏng. Sóng dọc: phương dao động trùng phương truyền, truyền trong rắn, lỏng và khí.
}
\end{ex}

% ===== Câu 87 =====
% ID: L11C2B9-22-TLN
% Mức: VD
\begin{ex}
% ID: L11C2B9-22-TLN
% Mức: VD
Từ kết quả của tình huống sau, nếu đại lượng cần tìm được nhân với hệ số 2 thì giá trị mới bằng bao nhiêu? Một sóng âm truyền trong không khí với tốc độ $340\,\text{m}$ /s trong $0,5\,\text{s}$. Quãng đường tính theo m bằng bao nhiêu?
\shortans{340}
\loigiai{
Xác định đúng dữ kiện và đơn vị, áp dụng hệ thức của dạng bài rồi tính được kết quả 340.
}
\end{ex}

% ===== Câu 88 =====
% ID: L11C2B9-22-TN
% Mức: VD
\dangbt{Nhận biết và vận dụng các đặc trưng của sóng âm}
\begin{ex}
% ID: L11C2B9-22-TN
% Mức: VD
Kết luận nào phù hợp nhất khi giải bài toán sau: Phân biệt độ cao và độ to của âm bằng hai đặc trưng vật lí tương ứng.
\choice
{Đại lượng cần tìm luôn bằng 0 trong mọi trường hợp.}
{Kết quả không phụ thuộc môi trường, tần số hoặc điều kiện biên.}
{\True Độ cao phụ thuộc chủ yếu vào tần số}
{Không cần sử dụng định nghĩa hay dữ kiện đã cho.}
\loigiai{
Đáp án C. Độ cao phụ thuộc chủ yếu vào tần số. Đối chiếu định nghĩa, hệ thức hoặc dữ kiện của bài toán cho kết luận trên.
}
\end{ex}

% ===== Câu 89 =====
% ID: L11C2B9-23-DS
% Mức: TH
\dangbt{Phân tích sự truyền năng lượng của sóng cơ}
\begin{ex}
% ID: L11C2B9-23-DS
% Mức: TH
Khi âm truyền từ loa đến tai.
\choiceTF
{\True Năng lượng được truyền qua không khí.}
{\True Các lớp không khí dao động quanh vị trí cân bằng.}
{\True Màng nhĩ dao động vì nhận năng lượng âm.}
{Luồng không khí từ loa phải đi hết tới tai.}
\loigiai{
\begin{itemchoice}
\item (Đúng) Năng lượng được truyền qua không khí. (Vì): Năng lượng được truyền qua không khí. b) Đ: Các lớp không khí dao động quanh vị trí cân bằng. c) Đ: Màng nhĩ dao động vì nhận năng lượng âm. d) S: Luồng không khí từ loa phải đi hết tới tai
\item (Đúng) Các lớp không khí dao động quanh vị trí cân bằng.
\item (Đúng) Màng nhĩ dao động vì nhận năng lượng âm.
\item (Sai) Luồng không khí từ loa phải đi hết tới tai.
\end{itemchoice}
}
\end{ex}

% ===== Câu 90 =====
% ID: L11C2B9-23-TL
% Mức: TH
\dangbt{Phân biệt sóng ngang, sóng dọc và môi trường truyền sóng}
\begin{ex}
% ID: L11C2B9-23-TL
% Mức: TH
Giải thích vì sao sóng âm trong không khí là sóng dọc.
\choiceTF

\loigiai{
Các lớp không khí dao động nén-dãn theo đúng phương truyền âm, tạo các vùng áp suất cao và thấp lan truyền.
}
\end{ex}

% ===== Câu 91 =====
% ID: L11C2B9-23-TLN
% Mức: VD
\begin{ex}
% ID: L11C2B9-23-TLN
% Mức: VD
Từ kết quả của tình huống sau, nếu đại lượng cần tìm được nhân với hệ số 0.5 thì giá trị mới bằng bao nhiêu? Sóng dọc có tần số $200\,\text{Hz}$ và tốc độ $300\,\text{m}$ /s. Bước sóng tính theo m bằng bao nhiêu?
\shortans{0.75}
\loigiai{
Xác định đúng dữ kiện và đơn vị, áp dụng hệ thức của dạng bài rồi tính được kết quả 0.75.
}
\end{ex}

% ===== Câu 92 =====
% ID: L11C2B9-23-TN
% Mức: VD
\dangbt{Nhận biết và vận dụng các đặc trưng của sóng âm}
\begin{ex}
% ID: L11C2B9-23-TN
% Mức: VD
Kết luận nào phù hợp nhất khi giải bài toán sau: Giải thích vì sao cùng một nốt nhạc truyền từ không khí vào nước vẫn có tần số không đổi nhưng bước sóng thay đổi.
\choice
{Không cần sử dụng định nghĩa hay dữ kiện đã cho.}
{Đại lượng cần tìm luôn bằng 0 trong mọi trường hợp.}
{Kết quả không phụ thuộc môi trường, tần số hoặc điều kiện biên.}
{\True Tần số do nguồn quyết định và liên tục qua mặt phân cách}
\loigiai{
Đáp án D. Tần số do nguồn quyết định và liên tục qua mặt phân cách. Đối chiếu định nghĩa, hệ thức hoặc dữ kiện của bài toán cho kết luận trên.
}
\end{ex}

% ===== Câu 93 =====
% ID: L11C2B9-24-DS
% Mức: VD
\dangbt{Phân tích sự truyền năng lượng của sóng cơ}
\begin{ex}
% ID: L11C2B9-24-DS
% Mức: VD
Xét các nhận định tổng hợp thuộc dạng “Phân tích sự truyền năng lượng của sóng cơ”.
\choiceTF
{\True dao động của các phần tử môi trường.}
{\True tăng.}
{\True năng lượng.}
{chuyển động tịnh tiến của toàn môi trường.}
\loigiai{
\begin{itemchoice}
\item (Đúng) dao động của các phần tử môi trường. (Vì): ao động của các phần tử môi trường. b) Đ: tăng. c) Đ: năng lượng. d) S: chuyển động tịnh tiến của toàn môi trường
\item (Đúng) tăng.
\item (Đúng) năng lượng.
\item (Sai) chuyển động tịnh tiến của toàn môi trường.
\end{itemchoice}
}
\end{ex}

% ===== Câu 94 =====
% ID: L11C2B9-24-TL
% Mức: VD
\dangbt{Phân biệt sóng ngang, sóng dọc và môi trường truyền sóng}
\begin{ex}
% ID: L11C2B9-24-TL
% Mức: VD
Nêu cách dùng hướng chuyển động của một phần tử để nhận biết loại sóng trên dây hoặc lò xo.
\choiceTF

\loigiai{
So sánh hướng dao động phần tử với hướng truyền: vuông góc là sóng ngang, song song/trùng phương là sóng dọc.
}
\end{ex}

% ===== Câu 95 =====
% ID: L11C2B9-24-TLN
% Mức: VD
\begin{ex}
% ID: L11C2B9-24-TLN
% Mức: VD
Từ kết quả của tình huống sau, nếu đại lượng cần tìm được nhân với hệ số 1.5 thì giá trị mới bằng bao nhiêu? Sóng trên dây truyền $12\,\text{m}$ trong $3\,\text{s}$. Tốc độ tính theo m/s bằng bao nhiêu?
\shortans{6}
\loigiai{
Xác định đúng dữ kiện và đơn vị, áp dụng hệ thức của dạng bài rồi tính được kết quả 6.
}
\end{ex}

% ===== Câu 96 =====
% ID: L11C2B9-24-TN
% Mức: VD
\dangbt{Nhận biết và vận dụng các đặc trưng của sóng âm}
\begin{ex}
% ID: L11C2B9-24-TN
% Mức: VD
Kết luận nào phù hợp nhất khi giải bài toán sau: Một âm $850\,\text{Hz}$ truyền trong không khí với tốc độ $340\,\text{m}$ /s. Tính bước sóng và cho biết âm thuộc vùng nghe được hay không.
\choice
{\True $\lambda=340/850=0,4$ m. Tần số $850\,\text{Hz}$ thuộc vùng âm nghe được của người.}
{Không cần sử dụng định nghĩa hay dữ kiện đã cho.}
{Đại lượng cần tìm luôn bằng 0 trong mọi trường hợp.}
{Kết quả không phụ thuộc môi trường, tần số hoặc điều kiện biên.}
\loigiai{
Đáp án A. $\lambda=340/850=0,4$ m. Tần số $850\,\text{Hz}$ thuộc vùng âm nghe được của người.. Đối chiếu định nghĩa, hệ thức hoặc dữ kiện của bài toán cho kết luận trên.
}
\end{ex}

% ===== Câu 97 =====
% ID: L11C2B9-25-DS
% Mức: VD
\dangbt{Phân tích sự truyền năng lượng của sóng cơ}
\begin{ex}
% ID: L11C2B9-25-DS
% Mức: VD
Đánh giá cách xử lí các dữ kiện định lượng của dạng “Phân tích sự truyền năng lượng của sóng cơ”.
\choiceTF
{\True Kết quả của bài toán thứ nhất là 4.}
{\True Kết quả của bài toán thứ hai là 5.}
{\True Kết quả của bài toán thứ ba là 24.}
{Có thể bỏ qua hoàn toàn đơn vị và điều kiện vật lí khi tính toán.}
\loigiai{
\begin{itemchoice}
\item (Đúng) Kết quả của bài toán thứ nhất là 4. (Vì): Kết quả của bài toán thứ nhất là 4. b) Đ: Kết quả của bài toán thứ hai là 5. c) Đ: Kết quả của bài toán thứ ba là 24. d) S: Có thể bỏ qua hoàn toàn đơn vị và điều kiện vật lí khi tính toán
\item (Đúng) Kết quả của bài toán thứ hai là 5.
\item (Đúng) Kết quả của bài toán thứ ba là 24.
\item (Sai) Có thể bỏ qua hoàn toàn đơn vị và điều kiện vật lí khi tính toán.
\end{itemchoice}
}
\end{ex}

% ===== Câu 98 =====
% ID: L11C2B9-25-TL
% Mức: VD
\dangbt{Phân biệt sóng ngang, sóng dọc và môi trường truyền sóng}
\begin{ex}
% ID: L11C2B9-25-TL
% Mức: VD
Phân tích các bước lập luận và chỉ ra điều kiện áp dụng trong bài toán: So sánh sóng ngang và sóng dọc theo phương dao động và môi trường truyền.
\choiceTF

\loigiai{
Sóng ngang: phương dao động vuông góc phương truyền, truyền trong chất rắn và trên bề mặt chất lỏng. Sóng dọc: phương dao động trùng phương truyền, truyền trong rắn, lỏng và khí. Cần nêu rõ điều kiện áp dụng, đổi đơn vị thống nhất và kiểm tra ý nghĩa vật lí của kết quả.
}
\end{ex}

% ===== Câu 99 =====
% ID: L11C2B9-25-TLN
% Mức: TH
\dangbt{Phân tích sự truyền năng lượng của sóng cơ}
\begin{ex}
% ID: L11C2B9-25-TLN
% Mức: TH
Nguồn truyền 24 $J$ năng lượng trong $6\,\text{s}$. Công suất truyền sóng tính theo $W$ bằng bao nhiêu?
\shortans{4}
\loigiai{
Xác định đúng dữ kiện và đơn vị, áp dụng hệ thức của dạng bài rồi tính được kết quả 4.
}
\end{ex}

% ===== Câu 100 =====
% ID: L11C2B9-25-TN
% Mức: VD
\dangbt{Nhận biết và vận dụng các đặc trưng của sóng âm}
\begin{ex}
% ID: L11C2B9-25-TN
% Mức: VD
Phương pháp phù hợp nhất cho dạng “Nhận biết và vận dụng các đặc trưng của sóng âm” là
\choice
{luôn coi mọi điểm dao động cùng pha}
{\True xác định dữ kiện, chọn đúng hệ thức hoặc định nghĩa, đổi đơn vị rồi kiểm tra kết quả}
{chỉ thay số mà không cần xét điều kiện áp dụng}
{bỏ qua đơn vị và chiều truyền sóng}
\loigiai{
Đáp án B. xác định dữ kiện, chọn đúng hệ thức hoặc định nghĩa, đổi đơn vị rồi kiểm tra kết quả. Đối chiếu định nghĩa, hệ thức hoặc dữ kiện của bài toán cho kết luận trên.
}
\end{ex}

% ===== Câu 101 =====
% ID: L11C2B9-26-DS
% Mức: TH
\dangbt{So sánh năng lượng và mức độ dao động của sóng}
\begin{ex}
% ID: L11C2B9-26-DS
% Mức: TH
Trong chuyên đề “So sánh năng lượng và mức độ dao động của sóng”, Một nguồn rung tạo sóng trên dây.
\choiceTF
{\True Nguồn thực hiện công lên dây.}
{\True Năng lượng truyền dọc theo dây.}
{\True Các phần tử dây nhận năng lượng và dao động.}
{Khối lượng dây được chuyển liên tục từ nguồn tới đầu kia.}
\loigiai{
\begin{itemchoice}
\item (Đúng) Nguồn thực hiện công lên dây. (Vì): Nguồn thực hiện công lên dây. b) Đ: Năng lượng truyền dọc theo dây. c) Đ: Các phần tử dây nhận năng lượng và dao động. d) S: Khối lượng dây được chuyển liên tục từ nguồn tới đầu kia
\item (Đúng) Năng lượng truyền dọc theo dây.
\item (Đúng) Các phần tử dây nhận năng lượng và dao động.
\item (Sai) Khối lượng dây được chuyển liên tục từ nguồn tới đầu kia.
\end{itemchoice}
}
\end{ex}

% ===== Câu 102 =====
% ID: L11C2B9-26-TL
% Mức: VD
\dangbt{Phân biệt sóng ngang, sóng dọc và môi trường truyền sóng}
\begin{ex}
% ID: L11C2B9-26-TL
% Mức: VD
Một học sinh giải bài sau nhưng bỏ qua điều kiện vật lí. Hãy sửa cách làm và trình bày lời giải đúng: Giải thích vì sao sóng âm trong không khí là sóng dọc.
\choiceTF

\loigiai{
Cách giải đúng: Các lớp không khí dao động nén-dãn theo đúng phương truyền âm, tạo các vùng áp suất cao và thấp lan truyền. Sai lầm cần tránh là dùng hệ thức khi chưa xác định điều kiện biên, môi trường hoặc đơn vị.
}
\end{ex}

% ===== Câu 103 =====
% ID: L11C2B9-26-TLN
% Mức: TH
\dangbt{Phân tích sự truyền năng lượng của sóng cơ}
\begin{ex}
% ID: L11C2B9-26-TLN
% Mức: TH
Một sóng mang 45 $J$ qua tiết diện trong $9\,\text{s}$. Công suất tính theo $W$ bằng bao nhiêu?
\shortans{5}
\loigiai{
Xác định đúng dữ kiện và đơn vị, áp dụng hệ thức của dạng bài rồi tính được kết quả 5.
}
\end{ex}

% ===== Câu 104 =====
% ID: L11C2B9-26-TN
% Mức: NB
\dangbt{Nhận dạng loại sóng trong các môi trường và thiết bị thực tế}
\begin{ex}
% ID: L11C2B9-26-TN
% Mức: NB
Trong chuyên đề “Nhận dạng loại sóng trong các môi trường và thiết bị thực tế”, Sóng trên một sợi dây căng có phương dao động vuông góc với dây là
\choice
{sóng âm trong khí}
{\True sóng ngang}
{sóng dọc}
{sóng điện từ}
\loigiai{
Đáp án B. sóng ngang. Đối chiếu định nghĩa, hệ thức hoặc dữ kiện của bài toán cho kết luận trên.
}
\end{ex}

% ===== Câu 105 =====
% ID: L11C2B9-27-DS
% Mức: TH
\dangbt{So sánh năng lượng và mức độ dao động của sóng}
\begin{ex}
% ID: L11C2B9-27-DS
% Mức: TH
Trong chuyên đề “So sánh năng lượng và mức độ dao động của sóng”, So sánh hai sóng cùng môi trường, cùng tần số nhưng biên độ khác nhau.
\choiceTF
{\True Sóng có biên độ lớn hơn thường mang nhiều năng lượng hơn.}
{\True Cả hai có cùng tốc độ truyền nếu môi trường không đổi.}
{\True Tần số do nguồn quyết định.}
{Biên độ không liên quan mức độ dao động của phần tử.}
\loigiai{
\begin{itemchoice}
\item (Đúng) Sóng có biên độ lớn hơn thường mang nhiều năng lượng hơn. (Vì): Sóng có biên độ lớn hơn thường mang nhiều năng lượng hơn. b) Đ: Cả hai có cùng tốc độ truyền nếu môi trường không đổi. c) Đ: Tần số do nguồn quyết định. d) S: Biên độ không liên quan mức độ dao động của phần tử
\item (Đúng) Cả hai có cùng tốc độ truyền nếu môi trường không đổi.
\item (Đúng) Tần số do nguồn quyết định.
\item (Sai) Biên độ không liên quan mức độ dao động của phần tử.
\end{itemchoice}
}
\end{ex}

% ===== Câu 106 =====
% ID: L11C2B9-27-TL
% Mức: VD
\dangbt{Phân biệt sóng ngang, sóng dọc và môi trường truyền sóng}
\begin{ex}
% ID: L11C2B9-27-TL
% Mức: VD
Mở rộng tình huống sau thành một ví dụ thực tế và trình bày phương pháp giải: Nêu cách dùng hướng chuyển động của một phần tử để nhận biết loại sóng trên dây hoặc lò xo.
\choiceTF

\loigiai{
So sánh hướng dao động phần tử với hướng truyền: vuông góc là sóng ngang, song song/trùng phương là sóng dọc. Có thể kiểm chứng bằng cách thay kết quả vào dữ kiện ban đầu và đối chiếu với hiện tượng thực tế.
}
\end{ex}

% ===== Câu 107 =====
% ID: L11C2B9-27-TLN
% Mức: TH
\dangbt{Phân tích sự truyền năng lượng của sóng cơ}
\begin{ex}
% ID: L11C2B9-27-TLN
% Mức: TH
Công suất sóng là $8\,\text{W}$. Năng lượng truyền trong $3\,\text{s}$ tính theo $J$ bằng bao nhiêu?
\shortans{24}
\loigiai{
Xác định đúng dữ kiện và đơn vị, áp dụng hệ thức của dạng bài rồi tính được kết quả 24.
}
\end{ex}

% ===== Câu 108 =====
% ID: L11C2B9-27-TN
% Mức: NB
\dangbt{Nhận dạng loại sóng trong các môi trường và thiết bị thực tế}
\begin{ex}
% ID: L11C2B9-27-TN
% Mức: NB
Trong chuyên đề “Nhận dạng loại sóng trong các môi trường và thiết bị thực tế”, Sóng âm truyền trong không khí là
\choice
{sóng đứng}
{sóng điện từ}
{\True sóng dọc}
{sóng ngang}
\loigiai{
Đáp án C. sóng dọc. Đối chiếu định nghĩa, hệ thức hoặc dữ kiện của bài toán cho kết luận trên.
}
\end{ex}

% ===== Câu 109 =====
% ID: L11C2B9-28-DS
% Mức: TH
\dangbt{So sánh năng lượng và mức độ dao động của sóng}
\begin{ex}
% ID: L11C2B9-28-DS
% Mức: TH
Trong chuyên đề “So sánh năng lượng và mức độ dao động của sóng”, Khi âm truyền từ loa đến tai.
\choiceTF
{\True Năng lượng được truyền qua không khí.}
{\True Các lớp không khí dao động quanh vị trí cân bằng.}
{\True Màng nhĩ dao động vì nhận năng lượng âm.}
{Luồng không khí từ loa phải đi hết tới tai.}
\loigiai{
\begin{itemchoice}
\item (Đúng) Năng lượng được truyền qua không khí. (Vì): Năng lượng được truyền qua không khí. b) Đ: Các lớp không khí dao động quanh vị trí cân bằng. c) Đ: Màng nhĩ dao động vì nhận năng lượng âm. d) S: Luồng không khí từ loa phải đi hết tới tai
\item (Đúng) Các lớp không khí dao động quanh vị trí cân bằng.
\item (Đúng) Màng nhĩ dao động vì nhận năng lượng âm.
\item (Sai) Luồng không khí từ loa phải đi hết tới tai.
\end{itemchoice}
}
\end{ex}

% ===== Câu 110 =====
% ID: L11C2B9-28-TL
% Mức: TH
\dangbt{Phân tích sự truyền năng lượng của sóng cơ}
\begin{ex}
% ID: L11C2B9-28-TL
% Mức: TH
Dùng quan điểm năng lượng để giải thích vì sao một vật nổi xa nguồn vẫn dao động khi sóng tới.
\choiceTF

\loigiai{
Sóng truyền năng lượng từ nguồn qua môi trường; khi tới vật nổi, năng lượng làm vật dao động dù vật chất của môi trường không đi từ nguồn tới vật.
}
\end{ex}

% ===== Câu 111 =====
% ID: L11C2B9-28-TLN
% Mức: VD
\begin{ex}
% ID: L11C2B9-28-TLN
% Mức: VD
Từ kết quả của tình huống sau, nếu đại lượng cần tìm được nhân với hệ số 2 thì giá trị mới bằng bao nhiêu? Nguồn truyền 24 $J$ năng lượng trong $6\,\text{s}$. Công suất truyền sóng tính theo $W$ bằng bao nhiêu?
\shortans{7.1}
\loigiai{
Xác định đúng dữ kiện và đơn vị, áp dụng hệ thức của dạng bài rồi tính được kết quả 8.
}
\end{ex}

% ===== Câu 112 =====
% ID: L11C2B9-28-TN
% Mức: TH
\dangbt{Nhận dạng loại sóng trong các môi trường và thiết bị thực tế}
\begin{ex}
% ID: L11C2B9-28-TN
% Mức: TH
Trong chuyên đề “Nhận dạng loại sóng trong các môi trường và thiết bị thực tế”, Sóng ngang cơ học truyền được trong
\choice
{mọi chất và chân không}
{chỉ chất khí}
{chỉ chất lỏng}
{\True chất rắn và trên bề mặt chất lỏng}
\loigiai{
Đáp án D. chất rắn và trên bề mặt chất lỏng. Đối chiếu định nghĩa, hệ thức hoặc dữ kiện của bài toán cho kết luận trên.
}
\end{ex}

% ===== Câu 113 =====
% ID: L11C2B9-29-DS
% Mức: VD
\dangbt{So sánh năng lượng và mức độ dao động của sóng}
\begin{ex}
% ID: L11C2B9-29-DS
% Mức: VD
Trong chuyên đề “So sánh năng lượng và mức độ dao động của sóng”, Xét các nhận định tổng hợp thuộc dạng “So sánh năng lượng và mức độ dao động của sóng”.
\choiceTF
{\True dao động của các phần tử môi trường.}
{\True tăng.}
{\True năng lượng.}
{chuyển động tịnh tiến của toàn môi trường.}
\loigiai{
\begin{itemchoice}
\item (Đúng) dao động của các phần tử môi trường. (Vì): ao động của các phần tử môi trường. b) Đ: tăng. c) Đ: năng lượng. d) S: chuyển động tịnh tiến của toàn môi trường
\item (Đúng) tăng.
\item (Đúng) năng lượng.
\item (Sai) chuyển động tịnh tiến của toàn môi trường.
\end{itemchoice}
}
\end{ex}

% ===== Câu 114 =====
% ID: L11C2B9-29-TL
% Mức: TH
\dangbt{Phân tích sự truyền năng lượng của sóng cơ}
\begin{ex}
% ID: L11C2B9-29-TL
% Mức: TH
Nêu sự khác nhau giữa tốc độ truyền năng lượng sóng và tốc độ dao động của phần tử môi trường.
\choiceTF

\loigiai{
Tốc độ truyền sóng mô tả sự lan truyền trạng thái và năng lượng; tốc độ phần tử là vận tốc dao động tại chỗ, biến thiên theo thời gian.
}
\end{ex}

% ===== Câu 115 =====
% ID: L11C2B9-29-TLN
% Mức: VD
\begin{ex}
% ID: L11C2B9-29-TLN
% Mức: VD
Từ kết quả của tình huống sau, nếu đại lượng cần tìm được nhân với hệ số 0.5 thì giá trị mới bằng bao nhiêu? Một sóng mang 45 $J$ qua tiết diện trong $9\,\text{s}$. Công suất tính theo $W$ bằng bao nhiêu?
\shortans{02/05/2026}
\loigiai{
Xác định đúng dữ kiện và đơn vị, áp dụng hệ thức của dạng bài rồi tính được kết quả 2.5.
}
\end{ex}

% ===== Câu 116 =====
% ID: L11C2B9-29-TN
% Mức: TH
\dangbt{Nhận dạng loại sóng trong các môi trường và thiết bị thực tế}
\begin{ex}
% ID: L11C2B9-29-TN
% Mức: TH
Trong chuyên đề “Nhận dạng loại sóng trong các môi trường và thiết bị thực tế”, Một sóng âm truyền trong không khí với tốc độ $340\,\text{m}$ /s trong $0,5\,\text{s}$. Quãng đường tính theo m bằng bao nhiêu? Chọn giá trị đúng.
\choice
{\True 170}
{340}
{85}
{171}
\loigiai{
Đáp án A. 170. Đối chiếu định nghĩa, hệ thức hoặc dữ kiện của bài toán cho kết luận trên.
}
\end{ex}

% ===== Câu 117 =====
% ID: L11C2B9-30-DS
% Mức: VD
\dangbt{So sánh năng lượng và mức độ dao động của sóng}
\begin{ex}
% ID: L11C2B9-30-DS
% Mức: VD
Trong chuyên đề “So sánh năng lượng và mức độ dao động của sóng”, Đánh giá cách xử lí các dữ kiện định lượng của dạng “So sánh năng lượng và mức độ dao động của sóng”.
\choiceTF
{\True Kết quả của bài toán thứ nhất là 4.}
{\True Kết quả của bài toán thứ hai là 5.}
{\True Kết quả của bài toán thứ ba là 24.}
{Có thể bỏ qua hoàn toàn đơn vị và điều kiện vật lí khi tính toán.}
\loigiai{
\begin{itemchoice}
\item (Đúng) Kết quả của bài toán thứ nhất là 4. (Vì): Kết quả của bài toán thứ nhất là 4. b) Đ: Kết quả của bài toán thứ hai là 5. c) Đ: Kết quả của bài toán thứ ba là 24. d) S: Có thể bỏ qua hoàn toàn đơn vị và điều kiện vật lí khi tính toán
\item (Đúng) Kết quả của bài toán thứ hai là 5.
\item (Đúng) Kết quả của bài toán thứ ba là 24.
\item (Sai) Có thể bỏ qua hoàn toàn đơn vị và điều kiện vật lí khi tính toán.
\end{itemchoice}
}
\end{ex}

% ===== Câu 118 =====
% ID: L11C2B9-30-TL
% Mức: VD
\dangbt{Phân tích sự truyền năng lượng của sóng cơ}
\begin{bt}
% ID: L11C2B9-30-TL
% Mức: VD
Một nguồn cấp 60 $J$ cho sóng trong $12\,\text{s}$ nhưng chỉ 48 $J$ đi qua tiết diện quan sát. Tính công suất tới và nêu nguyên nhân hao hụt.
\loigiai{
Công suất tới là $48/12=4$ W. Phần năng lượng còn lại có thể bị hấp thụ, tán xạ hoặc phản xạ.
}
\end{bt}

% ===== Câu 119 =====
% ID: L11C2B9-30-TLN
% Mức: VD
\begin{ex}
% ID: L11C2B9-30-TLN
% Mức: VD
Từ kết quả của tình huống sau, nếu đại lượng cần tìm được nhân với hệ số 1.5 thì giá trị mới bằng bao nhiêu? Công suất sóng là $8\,\text{W}$. Năng lượng truyền trong $3\,\text{s}$ tính theo $J$ bằng bao nhiêu?
\shortans{36}
\loigiai{
Xác định đúng dữ kiện và đơn vị, áp dụng hệ thức của dạng bài rồi tính được kết quả 36.
}
\end{ex}

% ===== Câu 120 =====
% ID: L11C2B9-30-TN
% Mức: TH
\dangbt{Nhận dạng loại sóng trong các môi trường và thiết bị thực tế}
\begin{ex}
% ID: L11C2B9-30-TN
% Mức: TH
Trong chuyên đề “Nhận dạng loại sóng trong các môi trường và thiết bị thực tế”, Sóng dọc có tần số $200\,\text{Hz}$ và tốc độ $300\,\text{m}$ /s. Bước sóng tính theo m bằng bao nhiêu? Chọn giá trị đúng.
\choice
{02/05/2026}
{\True 01/05/2026}
{3}
{0.75}
\loigiai{
Đáp án B. 1.5. Đối chiếu định nghĩa, hệ thức hoặc dữ kiện của bài toán cho kết luận trên.
}
\end{ex}

% ===== Câu 121 =====
% ID: L11C2B9-31-TL
% Mức: VD
\dangbt{Phân tích sự truyền năng lượng của sóng cơ}
\begin{ex}
% ID: L11C2B9-31-TL
% Mức: VD
Phân tích các bước lập luận và chỉ ra điều kiện áp dụng trong bài toán: Dùng quan điểm năng lượng để giải thích vì sao một vật nổi xa nguồn vẫn dao động khi sóng tới.
\choiceTF

\loigiai{
Sóng truyền năng lượng từ nguồn qua môi trường; khi tới vật nổi, năng lượng làm vật dao động dù vật chất của môi trường không đi từ nguồn tới vật. Cần nêu rõ điều kiện áp dụng, đổi đơn vị thống nhất và kiểm tra ý nghĩa vật lí của kết quả.
}
\end{ex}

% ===== Câu 122 =====
% ID: L11C2B9-31-TLN
% Mức: TH
\dangbt{So sánh năng lượng và mức độ dao động của sóng}
\begin{ex}
% ID: L11C2B9-31-TLN
% Mức: TH
Trong chuyên đề “So sánh năng lượng và mức độ dao động của sóng”, Nguồn truyền 24 $J$ năng lượng trong $6\,\text{s}$. Công suất truyền sóng tính theo $W$ bằng bao nhiêu?
\shortans{4}
\loigiai{
Xác định đúng dữ kiện và đơn vị, áp dụng hệ thức của dạng bài rồi tính được kết quả 4.
}
\end{ex}

% ===== Câu 123 =====
% ID: L11C2B9-31-TN
% Mức: TH
\dangbt{Nhận dạng loại sóng trong các môi trường và thiết bị thực tế}
\begin{ex}
% ID: L11C2B9-31-TN
% Mức: TH
Trong chuyên đề “Nhận dạng loại sóng trong các môi trường và thiết bị thực tế”, Sóng trên dây truyền $12\,\text{m}$ trong $3\,\text{s}$. Tốc độ tính theo m/s bằng bao nhiêu? Chọn giá trị đúng.
\choice
{2}
{5}
{\True 4}
{8}
\loigiai{
Đáp án C. 4. Đối chiếu định nghĩa, hệ thức hoặc dữ kiện của bài toán cho kết luận trên.
}
\end{ex}

% ===== Câu 124 =====
% ID: L11C2B9-32-TL
% Mức: VD
\dangbt{Phân tích sự truyền năng lượng của sóng cơ}
\begin{ex}
% ID: L11C2B9-32-TL
% Mức: VD
Một học sinh giải bài sau nhưng bỏ qua điều kiện vật lí. Hãy sửa cách làm và trình bày lời giải đúng: Nêu sự khác nhau giữa tốc độ truyền năng lượng sóng và tốc độ dao động của phần tử môi trường.
\choiceTF

\loigiai{
Cách giải đúng: Tốc độ truyền sóng mô tả sự lan truyền trạng thái và năng lượng; tốc độ phần tử là vận tốc dao động tại chỗ, biến thiên theo thời gian. Sai lầm cần tránh là dùng hệ thức khi chưa xác định điều kiện biên, môi trường hoặc đơn vị.
}
\end{ex}

% ===== Câu 125 =====
% ID: L11C2B9-32-TLN
% Mức: TH
\dangbt{So sánh năng lượng và mức độ dao động của sóng}
\begin{ex}
% ID: L11C2B9-32-TLN
% Mức: TH
Trong chuyên đề “So sánh năng lượng và mức độ dao động của sóng”, Một sóng mang 45 $J$ qua tiết diện trong $9\,\text{s}$. Công suất tính theo $W$ bằng bao nhiêu?
\shortans{5}
\loigiai{
Xác định đúng dữ kiện và đơn vị, áp dụng hệ thức của dạng bài rồi tính được kết quả 5.
}
\end{ex}

% ===== Câu 126 =====
% ID: L11C2B9-32-TN
% Mức: VD
\dangbt{Nhận dạng loại sóng trong các môi trường và thiết bị thực tế}
\begin{ex}
% ID: L11C2B9-32-TN
% Mức: VD
Trong chuyên đề “Nhận dạng loại sóng trong các môi trường và thiết bị thực tế”, Kết luận nào phù hợp nhất khi giải bài toán sau: So sánh sóng ngang và sóng dọc theo phương dao động và môi trường truyền.
\choice
{Không cần sử dụng định nghĩa hay dữ kiện đã cho.}
{Đại lượng cần tìm luôn bằng 0 trong mọi trường hợp.}
{Kết quả không phụ thuộc môi trường, tần số hoặc điều kiện biên.}
{\True Sóng ngang: phương dao động vuông góc phương truyền, truyền trong chất rắn và trên bề mặt chất lỏng. Sóng dọc: phương dao động trùng phương truyền, truyền trong rắn, lỏng và khí.}
\loigiai{
Đáp án D. Sóng ngang: phương dao động vuông góc phương truyền, truyền trong chất rắn và trên bề mặt chất lỏng. Sóng dọc: phương dao động trùng phương truyền, truyền trong rắn, lỏng và khí.. Đối chiếu định nghĩa, hệ thức hoặc dữ kiện của bài toán cho kết luận trên.
}
\end{ex}

% ===== Câu 127 =====
% ID: L11C2B9-33-TL
% Mức: VD
\dangbt{Phân tích sự truyền năng lượng của sóng cơ}
\begin{ex}
% ID: L11C2B9-33-TL
% Mức: VD
Mở rộng tình huống sau thành một ví dụ thực tế và trình bày phương pháp giải: Một nguồn cấp 60 $J$ cho sóng trong $12\,\text{s}$ nhưng chỉ 48 $J$ đi qua tiết diện quan sát. Tính công suất tới và nêu nguyên nhân hao hụt.
\choiceTF

\loigiai{
Công suất tới là $48/12=4$ W. Phần năng lượng còn lại có thể bị hấp thụ, tán xạ hoặc phản xạ. Có thể kiểm chứng bằng cách thay kết quả vào dữ kiện ban đầu và đối chiếu với hiện tượng thực tế.
}
\end{ex}

% ===== Câu 128 =====
% ID: L11C2B9-33-TLN
% Mức: TH
\dangbt{So sánh năng lượng và mức độ dao động của sóng}
\begin{ex}
% ID: L11C2B9-33-TLN
% Mức: TH
Trong chuyên đề “So sánh năng lượng và mức độ dao động của sóng”, Công suất sóng là $8\,\text{W}$. Năng lượng truyền trong $3\,\text{s}$ tính theo $J$ bằng bao nhiêu?
\shortans{24}
\loigiai{
Xác định đúng dữ kiện và đơn vị, áp dụng hệ thức của dạng bài rồi tính được kết quả 24.
}
\end{ex}

% ===== Câu 129 =====
% ID: L11C2B9-33-TN
% Mức: VD
\dangbt{Nhận dạng loại sóng trong các môi trường và thiết bị thực tế}
\begin{ex}
% ID: L11C2B9-33-TN
% Mức: VD
Trong chuyên đề “Nhận dạng loại sóng trong các môi trường và thiết bị thực tế”, Kết luận nào phù hợp nhất khi giải bài toán sau: Giải thích vì sao sóng âm trong không khí là sóng dọc.
\choice
{\True Các lớp không khí dao động nén-dãn theo đúng phương truyền âm, tạo các vùng áp suất cao và thấp lan truyền.}
{Không cần sử dụng định nghĩa hay dữ kiện đã cho.}
{Đại lượng cần tìm luôn bằng 0 trong mọi trường hợp.}
{Kết quả không phụ thuộc môi trường, tần số hoặc điều kiện biên.}
\loigiai{
Đáp án A. Các lớp không khí dao động nén-dãn theo đúng phương truyền âm, tạo các vùng áp suất cao và thấp lan truyền.. Đối chiếu định nghĩa, hệ thức hoặc dữ kiện của bài toán cho kết luận trên.
}
\end{ex}

% ===== Câu 130 =====
% ID: L11C2B9-34-TL
% Mức: TH
\dangbt{So sánh năng lượng và mức độ dao động của sóng}
\begin{ex}
% ID: L11C2B9-34-TL
% Mức: TH
Trong chuyên đề “So sánh năng lượng và mức độ dao động của sóng”, Dùng quan điểm năng lượng để giải thích vì sao một vật nổi xa nguồn vẫn dao động khi sóng tới.
\choiceTF

\loigiai{
Sóng truyền năng lượng từ nguồn qua môi trường; khi tới vật nổi, năng lượng làm vật dao động dù vật chất của môi trường không đi từ nguồn tới vật.
}
\end{ex}

% ===== Câu 131 =====
% ID: L11C2B9-34-TLN
% Mức: VD
\begin{ex}
% ID: L11C2B9-34-TLN
% Mức: VD
Trong chuyên đề “So sánh năng lượng và mức độ dao động của sóng”, Từ kết quả của tình huống sau, nếu đại lượng cần tìm được nhân với hệ số 2 thì giá trị mới bằng bao nhiêu? Nguồn truyền 24 $J$ năng lượng trong $6\,\text{s}$. Công suất truyền sóng tính theo $W$ bằng bao nhiêu?
\shortans{8}
\loigiai{
Xác định đúng dữ kiện và đơn vị, áp dụng hệ thức của dạng bài rồi tính được kết quả 8.
}
\end{ex}

% ===== Câu 132 =====
% ID: L11C2B9-34-TN
% Mức: VD
\dangbt{Nhận dạng loại sóng trong các môi trường và thiết bị thực tế}
\begin{ex}
% ID: L11C2B9-34-TN
% Mức: VD
Trong chuyên đề “Nhận dạng loại sóng trong các môi trường và thiết bị thực tế”, Kết luận nào phù hợp nhất khi giải bài toán sau: Nêu cách dùng hướng chuyển động của một phần tử để nhận biết loại sóng trên dây hoặc lò xo.
\choice
{Kết quả không phụ thuộc môi trường, tần số hoặc điều kiện biên.}
{\True So sánh hướng dao động phần tử với hướng truyền: vuông góc là sóng ngang, song song/trùng phương là sóng dọc.}
{Không cần sử dụng định nghĩa hay dữ kiện đã cho.}
{Đại lượng cần tìm luôn bằng 0 trong mọi trường hợp.}
\loigiai{
Đáp án B. So sánh hướng dao động phần tử với hướng truyền: vuông góc là sóng ngang, song song/trùng phương là sóng dọc.. Đối chiếu định nghĩa, hệ thức hoặc dữ kiện của bài toán cho kết luận trên.
}
\end{ex}

% ===== Câu 133 =====
% ID: L11C2B9-35-TL
% Mức: TH
\dangbt{So sánh năng lượng và mức độ dao động của sóng}
\begin{ex}
% ID: L11C2B9-35-TL
% Mức: TH
Trong chuyên đề “So sánh năng lượng và mức độ dao động của sóng”, Nêu sự khác nhau giữa tốc độ truyền năng lượng sóng và tốc độ dao động của phần tử môi trường.
\choiceTF

\loigiai{
Tốc độ truyền sóng mô tả sự lan truyền trạng thái và năng lượng; tốc độ phần tử là vận tốc dao động tại chỗ, biến thiên theo thời gian.
}
\end{ex}

% ===== Câu 134 =====
% ID: L11C2B9-35-TLN
% Mức: VD
\begin{ex}
% ID: L11C2B9-35-TLN
% Mức: VD
Trong chuyên đề “So sánh năng lượng và mức độ dao động của sóng”, Từ kết quả của tình huống sau, nếu đại lượng cần tìm được nhân với hệ số 0.5 thì giá trị mới bằng bao nhiêu? Một sóng mang 45 $J$ qua tiết diện trong $9\,\text{s}$. Công suất tính theo $W$ bằng bao nhiêu?
\shortans{02/05/2026}
\loigiai{
Xác định đúng dữ kiện và đơn vị, áp dụng hệ thức của dạng bài rồi tính được kết quả 2.5.
}
\end{ex}

% ===== Câu 135 =====
% ID: L11C2B9-35-TN
% Mức: VD
\dangbt{Nhận dạng loại sóng trong các môi trường và thiết bị thực tế}
\begin{ex}
% ID: L11C2B9-35-TN
% Mức: VD
Trong chuyên đề “Nhận dạng loại sóng trong các môi trường và thiết bị thực tế”, Phương pháp phù hợp nhất cho dạng “Nhận dạng loại sóng trong các môi trường và thiết bị thực tế” là
\choice
{bỏ qua đơn vị và chiều truyền sóng}
{luôn coi mọi điểm dao động cùng pha}
{\True xác định dữ kiện, chọn đúng hệ thức hoặc định nghĩa, đổi đơn vị rồi kiểm tra kết quả}
{chỉ thay số mà không cần xét điều kiện áp dụng}
\loigiai{
Đáp án C. xác định dữ kiện, chọn đúng hệ thức hoặc định nghĩa, đổi đơn vị rồi kiểm tra kết quả. Đối chiếu định nghĩa, hệ thức hoặc dữ kiện của bài toán cho kết luận trên.
}
\end{ex}

% ===== Câu 136 =====
% ID: L11C2B9-36-TL
% Mức: VD
\dangbt{So sánh năng lượng và mức độ dao động của sóng}
\begin{bt}
% ID: L11C2B9-36-TL
% Mức: VD
Trong chuyên đề “So sánh năng lượng và mức độ dao động của sóng”, Một nguồn cấp 60 $J$ cho sóng trong $12\,\text{s}$ nhưng chỉ 48 $J$ đi qua tiết diện quan sát. Tính công suất tới và nêu nguyên nhân hao hụt.
\loigiai{
Công suất tới là $48/12=4$ W. Phần năng lượng còn lại có thể bị hấp thụ, tán xạ hoặc phản xạ.
}
\end{bt}

% ===== Câu 137 =====
% ID: L11C2B9-36-TLN
% Mức: VD
\begin{ex}
% ID: L11C2B9-36-TLN
% Mức: VD
Trong chuyên đề “So sánh năng lượng và mức độ dao động của sóng”, Từ kết quả của tình huống sau, nếu đại lượng cần tìm được nhân với hệ số 1.5 thì giá trị mới bằng bao nhiêu? Công suất sóng là $8\,\text{W}$. Năng lượng truyền trong $3\,\text{s}$ tính theo $J$ bằng bao nhiêu?
\shortans{36}
\loigiai{
Xác định đúng dữ kiện và đơn vị, áp dụng hệ thức của dạng bài rồi tính được kết quả 36.
}
\end{ex}

% ===== Câu 138 =====
% ID: L11C2B9-36-TN
% Mức: NB
\dangbt{Phân biệt sóng ngang, sóng dọc và môi trường truyền sóng}
\begin{ex}
% ID: L11C2B9-36-TN
% Mức: NB
Sóng trên một sợi dây căng có phương dao động vuông góc với dây là
\choice
{sóng điện từ}
{sóng âm trong khí}
{\True sóng ngang}
{sóng dọc}
\loigiai{
Đáp án C. sóng ngang. Đối chiếu định nghĩa, hệ thức hoặc dữ kiện của bài toán cho kết luận trên.
}
\end{ex}

% ===== Câu 139 =====
% ID: L11C2B9-37-TL
% Mức: VD
\dangbt{So sánh năng lượng và mức độ dao động của sóng}
\begin{ex}
% ID: L11C2B9-37-TL
% Mức: VD
Trong chuyên đề “So sánh năng lượng và mức độ dao động của sóng”, Phân tích các bước lập luận và chỉ ra điều kiện áp dụng trong bài toán: Dùng quan điểm năng lượng để giải thích vì sao một vật nổi xa nguồn vẫn dao động khi sóng tới.
\choiceTF

\loigiai{
Sóng truyền năng lượng từ nguồn qua môi trường; khi tới vật nổi, năng lượng làm vật dao động dù vật chất của môi trường không đi từ nguồn tới vật. Cần nêu rõ điều kiện áp dụng, đổi đơn vị thống nhất và kiểm tra ý nghĩa vật lí của kết quả.
}
\end{ex}

% ===== Câu 140 =====
% ID: L11C2B9-37-TN
% Mức: NB
\dangbt{Phân biệt sóng ngang, sóng dọc và môi trường truyền sóng}
\begin{ex}
% ID: L11C2B9-37-TN
% Mức: NB
Sóng âm truyền trong không khí là
\choice
{sóng ngang}
{sóng đứng}
{sóng điện từ}
{\True sóng dọc}
\loigiai{
Đáp án D. sóng dọc. Đối chiếu định nghĩa, hệ thức hoặc dữ kiện của bài toán cho kết luận trên.
}
\end{ex}

% ===== Câu 141 =====
% ID: L11C2B9-38-TL
% Mức: VD
\dangbt{So sánh năng lượng và mức độ dao động của sóng}
\begin{ex}
% ID: L11C2B9-38-TL
% Mức: VD
Trong chuyên đề “So sánh năng lượng và mức độ dao động của sóng”, Một học sinh giải bài sau nhưng bỏ qua điều kiện vật lí. Hãy sửa cách làm và trình bày lời giải đúng: Nêu sự khác nhau giữa tốc độ truyền năng lượng sóng và tốc độ dao động của phần tử môi trường.
\choiceTF

\loigiai{
Cách giải đúng: Tốc độ truyền sóng mô tả sự lan truyền trạng thái và năng lượng; tốc độ phần tử là vận tốc dao động tại chỗ, biến thiên theo thời gian. Sai lầm cần tránh là dùng hệ thức khi chưa xác định điều kiện biên, môi trường hoặc đơn vị.
}
\end{ex}

% ===== Câu 142 =====
% ID: L11C2B9-38-TN
% Mức: TH
\dangbt{Phân biệt sóng ngang, sóng dọc và môi trường truyền sóng}
\begin{ex}
% ID: L11C2B9-38-TN
% Mức: TH
Sóng ngang cơ học truyền được trong
\choice
{\True chất rắn và trên bề mặt chất lỏng}
{mọi chất và chân không}
{chỉ chất khí}
{chỉ chất lỏng}
\loigiai{
Đáp án A. chất rắn và trên bề mặt chất lỏng. Đối chiếu định nghĩa, hệ thức hoặc dữ kiện của bài toán cho kết luận trên.
}
\end{ex}

% ===== Câu 143 =====
% ID: L11C2B9-39-TL
% Mức: VD
\dangbt{So sánh năng lượng và mức độ dao động của sóng}
\begin{ex}
% ID: L11C2B9-39-TL
% Mức: VD
Trong chuyên đề “So sánh năng lượng và mức độ dao động của sóng”, Mở rộng tình huống sau thành một ví dụ thực tế và trình bày phương pháp giải: Một nguồn cấp 60 $J$ cho sóng trong $12\,\text{s}$ nhưng chỉ 48 $J$ đi qua tiết diện quan sát. Tính công suất tới và nêu nguyên nhân hao hụt.
\choiceTF

\loigiai{
Công suất tới là $48/12=4$ W. Phần năng lượng còn lại có thể bị hấp thụ, tán xạ hoặc phản xạ. Có thể kiểm chứng bằng cách thay kết quả vào dữ kiện ban đầu và đối chiếu với hiện tượng thực tế.
}
\end{ex}

% ===== Câu 144 =====
% ID: L11C2B9-39-TN
% Mức: TH
\dangbt{Phân biệt sóng ngang, sóng dọc và môi trường truyền sóng}
\begin{ex}
% ID: L11C2B9-39-TN
% Mức: TH
Một sóng âm truyền trong không khí với tốc độ $340\,\text{m}$ /s trong $0,5\,\text{s}$. Quãng đường tính theo m bằng bao nhiêu? Chọn giá trị đúng.
\choice
{171}
{\True 170}
{340}
{85}
\loigiai{
Đáp án B. 170. Đối chiếu định nghĩa, hệ thức hoặc dữ kiện của bài toán cho kết luận trên.
}
\end{ex}

% ===== Câu 145 =====
% ID: L11C2B9-40-TN
% Mức: TH
\begin{ex}
% ID: L11C2B9-40-TN
% Mức: TH
Sóng dọc có tần số $200\,\text{Hz}$ và tốc độ $300\,\text{m}$ /s. Bước sóng tính theo m bằng bao nhiêu? Chọn giá trị đúng.
\choice
{0.75}
{02/05/2026}
{\True 01/05/2026}
{3}
\loigiai{
Đáp án C. 1.5. Đối chiếu định nghĩa, hệ thức hoặc dữ kiện của bài toán cho kết luận trên.
}
\end{ex}

% ===== Câu 146 =====
% ID: L11C2B9-41-TN
% Mức: TH
\begin{ex}
% ID: L11C2B9-41-TN
% Mức: TH
Sóng trên dây truyền $12\,\text{m}$ trong $3\,\text{s}$. Tốc độ tính theo m/s bằng bao nhiêu? Chọn giá trị đúng.
\choice
{8}
{2}
{5}
{\True 4}
\loigiai{
Đáp án D. 4. Đối chiếu định nghĩa, hệ thức hoặc dữ kiện của bài toán cho kết luận trên.
}
\end{ex}

% ===== Câu 147 =====
% ID: L11C2B9-42-TN
% Mức: VD
\begin{ex}
% ID: L11C2B9-42-TN
% Mức: VD
Kết luận nào phù hợp nhất khi giải bài toán sau: So sánh sóng ngang và sóng dọc theo phương dao động và môi trường truyền.
\choice
{\True Sóng ngang: phương dao động vuông góc phương truyền, truyền trong chất rắn và trên bề mặt chất lỏng. Sóng dọc: phương dao động trùng phương truyền, truyền trong rắn, lỏng và khí.}
{Không cần sử dụng định nghĩa hay dữ kiện đã cho.}
{Đại lượng cần tìm luôn bằng 0 trong mọi trường hợp.}
{Kết quả không phụ thuộc môi trường, tần số hoặc điều kiện biên.}
\loigiai{
Đáp án A. Sóng ngang: phương dao động vuông góc phương truyền, truyền trong chất rắn và trên bề mặt chất lỏng. Sóng dọc: phương dao động trùng phương truyền, truyền trong rắn, lỏng và khí.. Đối chiếu định nghĩa, hệ thức hoặc dữ kiện của bài toán cho kết luận trên.
}
\end{ex}

% ===== Câu 148 =====
% ID: L11C2B9-43-TN
% Mức: VD
\begin{ex}
% ID: L11C2B9-43-TN
% Mức: VD
Kết luận nào phù hợp nhất khi giải bài toán sau: Giải thích vì sao sóng âm trong không khí là sóng dọc.
\choice
{Kết quả không phụ thuộc môi trường, tần số hoặc điều kiện biên.}
{\True Các lớp không khí dao động nén-dãn theo đúng phương truyền âm, tạo các vùng áp suất cao và thấp lan truyền.}
{Không cần sử dụng định nghĩa hay dữ kiện đã cho.}
{Đại lượng cần tìm luôn bằng 0 trong mọi trường hợp.}
\loigiai{
Đáp án B. Các lớp không khí dao động nén-dãn theo đúng phương truyền âm, tạo các vùng áp suất cao và thấp lan truyền.. Đối chiếu định nghĩa, hệ thức hoặc dữ kiện của bài toán cho kết luận trên.
}
\end{ex}

% ===== Câu 149 =====
% ID: L11C2B9-44-TN
% Mức: VD
\begin{ex}
% ID: L11C2B9-44-TN
% Mức: VD
Kết luận nào phù hợp nhất khi giải bài toán sau: Nêu cách dùng hướng chuyển động của một phần tử để nhận biết loại sóng trên dây hoặc lò xo.
\choice
{Đại lượng cần tìm luôn bằng 0 trong mọi trường hợp.}
{Kết quả không phụ thuộc môi trường, tần số hoặc điều kiện biên.}
{\True So sánh hướng dao động phần tử với hướng truyền: vuông góc là sóng ngang, song song/trùng phương là sóng dọc.}
{Không cần sử dụng định nghĩa hay dữ kiện đã cho.}
\loigiai{
Đáp án C. So sánh hướng dao động phần tử với hướng truyền: vuông góc là sóng ngang, song song/trùng phương là sóng dọc.. Đối chiếu định nghĩa, hệ thức hoặc dữ kiện của bài toán cho kết luận trên.
}
\end{ex}

% ===== Câu 150 =====
% ID: L11C2B9-45-TN
% Mức: VD
\begin{ex}
% ID: L11C2B9-45-TN
% Mức: VD
Phương pháp phù hợp nhất cho dạng “Phân biệt sóng ngang, sóng dọc và môi trường truyền sóng” là
\choice
{chỉ thay số mà không cần xét điều kiện áp dụng}
{bỏ qua đơn vị và chiều truyền sóng}
{luôn coi mọi điểm dao động cùng pha}
{\True xác định dữ kiện, chọn đúng hệ thức hoặc định nghĩa, đổi đơn vị rồi kiểm tra kết quả}
\loigiai{
Đáp án D. xác định dữ kiện, chọn đúng hệ thức hoặc định nghĩa, đổi đơn vị rồi kiểm tra kết quả. Đối chiếu định nghĩa, hệ thức hoặc dữ kiện của bài toán cho kết luận trên.
}
\end{ex}

% ===== Câu 151 =====
% ID: L11C2B9-46-TN
% Mức: NB
\dangbt{Phân tích sự truyền năng lượng của sóng cơ}
\begin{ex}
% ID: L11C2B9-46-TN
% Mức: NB
Năng lượng sóng cơ tại một vùng môi trường là năng lượng
\choice
{chuyển động tịnh tiến của toàn môi trường}
{điện trường tĩnh}
{hóa học của nguồn}
{\True dao động của các phần tử môi trường}
\loigiai{
Đáp án D. dao động của các phần tử môi trường. Đối chiếu định nghĩa, hệ thức hoặc dữ kiện của bài toán cho kết luận trên.
}
\end{ex}

% ===== Câu 152 =====
% ID: L11C2B9-47-TN
% Mức: NB
\begin{ex}
% ID: L11C2B9-47-TN
% Mức: NB
Khi biên độ sóng tăng, năng lượng sóng thông thường
\choice
{\True tăng}
{giảm về 0}
{không đổi}
{đổi dấu}
\loigiai{
Đáp án A. tăng. Đối chiếu định nghĩa, hệ thức hoặc dữ kiện của bài toán cho kết luận trên.
}
\end{ex}

% ===== Câu 153 =====
% ID: L11C2B9-48-TN
% Mức: TH
\begin{ex}
% ID: L11C2B9-48-TN
% Mức: TH
Một phao dao động mạnh hơn khi sóng lớn tới cho thấy sóng đã truyền
\choice
{nhiệt độ tuyệt đối}
{\True năng lượng}
{khối lượng nước}
{điện tích}
\loigiai{
Đáp án B. năng lượng. Đối chiếu định nghĩa, hệ thức hoặc dữ kiện của bài toán cho kết luận trên.
}
\end{ex}

% ===== Câu 154 =====
% ID: L11C2B9-49-TN
% Mức: TH
\begin{ex}
% ID: L11C2B9-49-TN
% Mức: TH
Nguồn truyền 24 $J$ năng lượng trong $6\,\text{s}$. Công suất truyền sóng tính theo $W$ bằng bao nhiêu? Chọn giá trị đúng.
\choice
{2}
{5}
{\True 4}
{8}
\loigiai{
Đáp án C. 4. Đối chiếu định nghĩa, hệ thức hoặc dữ kiện của bài toán cho kết luận trên.
}
\end{ex}

% ===== Câu 155 =====
% ID: L11C2B9-50-TN
% Mức: TH
\begin{ex}
% ID: L11C2B9-50-TN
% Mức: TH
Một sóng mang 45 $J$ qua tiết diện trong $9\,\text{s}$. Công suất tính theo $W$ bằng bao nhiêu? Chọn giá trị đúng.
\choice
{10}
{02/05/2026}
{6}
{\True 5}
\loigiai{
Đáp án D. 5. Đối chiếu định nghĩa, hệ thức hoặc dữ kiện của bài toán cho kết luận trên.
}
\end{ex}

% ===== Câu 156 =====
% ID: L11C2B9-51-TN
% Mức: TH
\begin{ex}
% ID: L11C2B9-51-TN
% Mức: TH
Công suất sóng là $8\,\text{W}$. Năng lượng truyền trong $3\,\text{s}$ tính theo $J$ bằng bao nhiêu? Chọn giá trị đúng.
\choice
{\True 24}
{48}
{12}
{25}
\loigiai{
Đáp án A. 24. Đối chiếu định nghĩa, hệ thức hoặc dữ kiện của bài toán cho kết luận trên.
}
\end{ex}

% ===== Câu 157 =====
% ID: L11C2B9-52-TN
% Mức: VD
\begin{ex}
% ID: L11C2B9-52-TN
% Mức: VD
Kết luận nào phù hợp nhất khi giải bài toán sau: Dùng quan điểm năng lượng để giải thích vì sao một vật nổi xa nguồn vẫn dao động khi sóng tới.
\choice
{Kết quả không phụ thuộc môi trường, tần số hoặc điều kiện biên.}
{\True Sóng truyền năng lượng từ nguồn qua môi trường}
{Không cần sử dụng định nghĩa hay dữ kiện đã cho.}
{Đại lượng cần tìm luôn bằng 0 trong mọi trường hợp.}
\loigiai{
Đáp án B. Sóng truyền năng lượng từ nguồn qua môi trường. Đối chiếu định nghĩa, hệ thức hoặc dữ kiện của bài toán cho kết luận trên.
}
\end{ex}

% ===== Câu 158 =====
% ID: L11C2B9-53-TN
% Mức: VD
\begin{ex}
% ID: L11C2B9-53-TN
% Mức: VD
Kết luận nào phù hợp nhất khi giải bài toán sau: Nêu sự khác nhau giữa tốc độ truyền năng lượng sóng và tốc độ dao động của phần tử môi trường.
\choice
{Đại lượng cần tìm luôn bằng 0 trong mọi trường hợp.}
{Kết quả không phụ thuộc môi trường, tần số hoặc điều kiện biên.}
{\True Tốc độ truyền sóng mô tả sự lan truyền trạng thái và năng lượng}
{Không cần sử dụng định nghĩa hay dữ kiện đã cho.}
\loigiai{
Đáp án C. Tốc độ truyền sóng mô tả sự lan truyền trạng thái và năng lượng. Đối chiếu định nghĩa, hệ thức hoặc dữ kiện của bài toán cho kết luận trên.
}
\end{ex}

% ===== Câu 159 =====
% ID: L11C2B9-54-TN
% Mức: VD
\begin{ex}
% ID: L11C2B9-54-TN
% Mức: VD
Kết luận nào phù hợp nhất khi giải bài toán sau: Một nguồn cấp 60 $J$ cho sóng trong $12\,\text{s}$ nhưng chỉ 48 $J$ đi qua tiết diện quan sát. Tính công suất tới và nêu nguyên nhân hao hụt.
\choice
{Không cần sử dụng định nghĩa hay dữ kiện đã cho.}
{Đại lượng cần tìm luôn bằng 0 trong mọi trường hợp.}
{Kết quả không phụ thuộc môi trường, tần số hoặc điều kiện biên.}
{\True Công suất tới là $48/12=4$ W. Phần năng lượng còn lại có thể bị hấp thụ, tán xạ hoặc phản xạ.}
\loigiai{
Đáp án D. Công suất tới là $48/12=4$ W. Phần năng lượng còn lại có thể bị hấp thụ, tán xạ hoặc phản xạ.. Đối chiếu định nghĩa, hệ thức hoặc dữ kiện của bài toán cho kết luận trên.
}
\end{ex}

% ===== Câu 160 =====
% ID: L11C2B9-55-TN
% Mức: VD
\begin{ex}
% ID: L11C2B9-55-TN
% Mức: VD
Phương pháp phù hợp nhất cho dạng “Phân tích sự truyền năng lượng của sóng cơ” là
\choice
{\True xác định dữ kiện, chọn đúng hệ thức hoặc định nghĩa, đổi đơn vị rồi kiểm tra kết quả}
{chỉ thay số mà không cần xét điều kiện áp dụng}
{bỏ qua đơn vị và chiều truyền sóng}
{luôn coi mọi điểm dao động cùng pha}
\loigiai{
Đáp án A. xác định dữ kiện, chọn đúng hệ thức hoặc định nghĩa, đổi đơn vị rồi kiểm tra kết quả. Đối chiếu định nghĩa, hệ thức hoặc dữ kiện của bài toán cho kết luận trên.
}
\end{ex}

% ===== Câu 161 =====
% ID: L11C2B9-56-TN
% Mức: NB
\dangbt{So sánh năng lượng và mức độ dao động của sóng}
\begin{ex}
% ID: L11C2B9-56-TN
% Mức: NB
Trong chuyên đề “So sánh năng lượng và mức độ dao động của sóng”, Năng lượng sóng cơ tại một vùng môi trường là năng lượng
\choice
{điện trường tĩnh}
{hóa học của nguồn}
{\True dao động của các phần tử môi trường}
{chuyển động tịnh tiến của toàn môi trường}
\loigiai{
Đáp án C. dao động của các phần tử môi trường. Đối chiếu định nghĩa, hệ thức hoặc dữ kiện của bài toán cho kết luận trên.
}
\end{ex}

% ===== Câu 162 =====
% ID: L11C2B9-57-TN
% Mức: NB
\begin{ex}
% ID: L11C2B9-57-TN
% Mức: NB
Trong chuyên đề “So sánh năng lượng và mức độ dao động của sóng”, Khi biên độ sóng tăng, năng lượng sóng thông thường
\choice
{giảm về 0}
{không đổi}
{đổi dấu}
{\True tăng}
\loigiai{
Đáp án D. tăng. Đối chiếu định nghĩa, hệ thức hoặc dữ kiện của bài toán cho kết luận trên.
}
\end{ex}

% ===== Câu 163 =====
% ID: L11C2B9-58-TN
% Mức: TH
\begin{ex}
% ID: L11C2B9-58-TN
% Mức: TH
Trong chuyên đề “So sánh năng lượng và mức độ dao động của sóng”, Một phao dao động mạnh hơn khi sóng lớn tới cho thấy sóng đã truyền
\choice
{\True năng lượng}
{khối lượng nước}
{điện tích}
{nhiệt độ tuyệt đối}
\loigiai{
Đáp án A. năng lượng. Đối chiếu định nghĩa, hệ thức hoặc dữ kiện của bài toán cho kết luận trên.
}
\end{ex}

% ===== Câu 164 =====
% ID: L11C2B9-59-TN
% Mức: TH
\begin{ex}
% ID: L11C2B9-59-TN
% Mức: TH
Trong chuyên đề “So sánh năng lượng và mức độ dao động của sóng”, Nguồn truyền 24 $J$ năng lượng trong $6\,\text{s}$. Công suất truyền sóng tính theo $W$ bằng bao nhiêu? Chọn giá trị đúng.
\choice
{5}
{\True 4}
{8}
{2}
\loigiai{
Đáp án B. 4. Đối chiếu định nghĩa, hệ thức hoặc dữ kiện của bài toán cho kết luận trên.
}
\end{ex}

% ===== Câu 165 =====
% ID: L11C2B9-60-TN
% Mức: TH
\begin{ex}
% ID: L11C2B9-60-TN
% Mức: TH
Trong chuyên đề “So sánh năng lượng và mức độ dao động của sóng”, Một sóng mang 45 $J$ qua tiết diện trong $9\,\text{s}$. Công suất tính theo $W$ bằng bao nhiêu? Chọn giá trị đúng.
\choice
{02/05/2026}
{6}
{\True 5}
{10}
\loigiai{
Đáp án C. 5. Đối chiếu định nghĩa, hệ thức hoặc dữ kiện của bài toán cho kết luận trên.
}
\end{ex}

% ===== Câu 166 =====
% ID: L11C2B9-61-TN
% Mức: TH
\begin{ex}
% ID: L11C2B9-61-TN
% Mức: TH
Trong chuyên đề “So sánh năng lượng và mức độ dao động của sóng”, Công suất sóng là $8\,\text{W}$. Năng lượng truyền trong $3\,\text{s}$ tính theo $J$ bằng bao nhiêu? Chọn giá trị đúng.
\choice
{48}
{12}
{25}
{\True 24}
\loigiai{
Đáp án D. 24. Đối chiếu định nghĩa, hệ thức hoặc dữ kiện của bài toán cho kết luận trên.
}
\end{ex}

% ===== Câu 167 =====
% ID: L11C2B9-62-TN
% Mức: VD
\begin{ex}
% ID: L11C2B9-62-TN
% Mức: VD
Trong chuyên đề “So sánh năng lượng và mức độ dao động của sóng”, Kết luận nào phù hợp nhất khi giải bài toán sau: Dùng quan điểm năng lượng để giải thích vì sao một vật nổi xa nguồn vẫn dao động khi sóng tới.
\choice
{\True Sóng truyền năng lượng từ nguồn qua môi trường}
{Không cần sử dụng định nghĩa hay dữ kiện đã cho.}
{Đại lượng cần tìm luôn bằng 0 trong mọi trường hợp.}
{Kết quả không phụ thuộc môi trường, tần số hoặc điều kiện biên.}
\loigiai{
Đáp án A. Sóng truyền năng lượng từ nguồn qua môi trường. Đối chiếu định nghĩa, hệ thức hoặc dữ kiện của bài toán cho kết luận trên.
}
\end{ex}

% ===== Câu 168 =====
% ID: L11C2B9-63-TN
% Mức: VD
\begin{ex}
% ID: L11C2B9-63-TN
% Mức: VD
Trong chuyên đề “So sánh năng lượng và mức độ dao động của sóng”, Kết luận nào phù hợp nhất khi giải bài toán sau: Nêu sự khác nhau giữa tốc độ truyền năng lượng sóng và tốc độ dao động của phần tử môi trường.
\choice
{Kết quả không phụ thuộc môi trường, tần số hoặc điều kiện biên.}
{\True Tốc độ truyền sóng mô tả sự lan truyền trạng thái và năng lượng}
{Không cần sử dụng định nghĩa hay dữ kiện đã cho.}
{Đại lượng cần tìm luôn bằng 0 trong mọi trường hợp.}
\loigiai{
Đáp án B. Tốc độ truyền sóng mô tả sự lan truyền trạng thái và năng lượng. Đối chiếu định nghĩa, hệ thức hoặc dữ kiện của bài toán cho kết luận trên.
}
\end{ex}

% ===== Câu 169 =====
% ID: L11C2B9-64-TN
% Mức: VD
\begin{ex}
% ID: L11C2B9-64-TN
% Mức: VD
Trong chuyên đề “So sánh năng lượng và mức độ dao động của sóng”, Kết luận nào phù hợp nhất khi giải bài toán sau: Một nguồn cấp 60 $J$ cho sóng trong $12\,\text{s}$ nhưng chỉ 48 $J$ đi qua tiết diện quan sát. Tính công suất tới và nêu nguyên nhân hao hụt.
\choice
{Đại lượng cần tìm luôn bằng 0 trong mọi trường hợp.}
{Kết quả không phụ thuộc môi trường, tần số hoặc điều kiện biên.}
{\True Công suất tới là $48/12=4$ W. Phần năng lượng còn lại có thể bị hấp thụ, tán xạ hoặc phản xạ.}
{Không cần sử dụng định nghĩa hay dữ kiện đã cho.}
\loigiai{
Đáp án C. Công suất tới là $48/12=4$ W. Phần năng lượng còn lại có thể bị hấp thụ, tán xạ hoặc phản xạ.. Đối chiếu định nghĩa, hệ thức hoặc dữ kiện của bài toán cho kết luận trên.
}
\end{ex}

% ===== Câu 170 =====
% ID: L11C2B9-65-TN
% Mức: VD
\begin{ex}
% ID: L11C2B9-65-TN
% Mức: VD
Trong chuyên đề “So sánh năng lượng và mức độ dao động của sóng”, Phương pháp phù hợp nhất cho dạng “So sánh năng lượng và mức độ dao động của sóng” là
\choice
{chỉ thay số mà không cần xét điều kiện áp dụng}
{bỏ qua đơn vị và chiều truyền sóng}
{luôn coi mọi điểm dao động cùng pha}
{\True xác định dữ kiện, chọn đúng hệ thức hoặc định nghĩa, đổi đơn vị rồi kiểm tra kết quả}
\loigiai{
Đáp án D. xác định dữ kiện, chọn đúng hệ thức hoặc định nghĩa, đổi đơn vị rồi kiểm tra kết quả. Đối chiếu định nghĩa, hệ thức hoặc dữ kiện của bài toán cho kết luận trên.
}
\end{ex}
